
\chapter{Numerical Methods}

\noindent
In this chapter we study numerical methods for solving a first order
differential equation
\[
y'=f(x,y).
\]


\medskip\noindent
SECTION~3.1 deals with
\href{http://www-history.mcs.st-and.ac.uk/Mathematicians/Euler.html}
{\color{blue}\itshape Euler's
  method}, which is
really too crude to be of much use in practical applications. However,
its simplicity allows for an introduction to the ideas required to
understand the better methods discussed in the other two sections.

\medskip\noindent
SECTION~3.2 discusses improvements on Euler's method.


\medskip\noindent
SECTION~3.3 deals with the
\href{http://www-history.mcs.st-and.ac.uk/Mathematicians/Runge.html}
{Runge\color{blue}}-\href{http://www-history.mcs.st-and.ac.uk/Mathematicians/Kutta.html}
{\color{blue}Kutta}
method, perhaps
the most widely used method for numerical solution of differential
equations.
 \vspace*{4in}

\centerline{95}

\newpage
\section{Euler's Method}

If an initial value problem
\begin{equation} \label{eq:3.1.1}
y'=f(x,y),\quad y(x_0)=y_0
 \end{equation}
can't be solved analytically,
 it's necessary to resort to
numerical methods to obtain useful  approximations to a solution of
\eqref{eq:3.1.1}. We'll consider such methods in this chapter.

We're interested in
computing approximate values of the solution of \eqref{eq:3.1.1} at
equally spaced points $x_0$, $x_1$, \dots, $x_n=b$ in an interval
$[x_0,b]$.
 Thus,
\[
x_i=x_0+ih,\quad i=0,1, \dots,n,
\]
where
\[
h={b-x_0\over n}.
\]
We'll denote the approximate values of the solution at these points
by $y_0$, $y_1$, \dots, $y_n$;   thus, $y_i$ is an approximation to
$y(x_i)$.
We'll call
\[
e_i=y(x_i)-y_i
\]
the \emph{\color{blue}error at the $i$th step}. Because of the initial
condition
$y(x_0)=y_0$, we'll always have $e_0=0$. However, in general
$e_i\ne0$ if $i>0$.

We encounter two sources of error
in applying a numerical method to solve an initial value problem:
\begin{itemize}
\item
The formulas defining the method are based on some sort of
approximation. Errors due to the inaccuracy of the approximation are
called \emph{\color{blue}truncation errors}.
\item
Computers do arithmetic with a fixed number of digits, and therefore
make errors in evaluating the formulas defining the numerical methods.
Errors due to the computer's inability to do exact arithmetic are
called \emph{\color{blue}roundoff errors}.
\end{itemize}


Since a careful analysis of roundoff error is beyond the scope of this
book, we'll consider only truncation errors.

\boxit{Euler's Method}


\noindent
The simplest numerical method for solving \eqref{eq:3.1.1} is
\href{http://www-history.mcs.st-and.ac.uk/Mathematicians/Euler.html}
{\color{blue}\it
Euler's method\/}.
This method is so crude that it
is seldom used in practice;     however, its simplicity makes it useful
for illustrative purposes.

Euler's method is based on the assumption that the tangent line to the
integral curve of \eqref{eq:3.1.1} at $(x_i,y(x_i))$ approximates the
integral curve over the interval $[x_i,x_{i+1}]$.
 Since the slope of the integral curve of
\eqref{eq:3.1.1} at $(x_i,y(x_i))$ is $y'(x_i)=f(x_i,y(x_i))$, the
equation of the tangent line to the integral curve at $(x_i,y(x_i))$
is
\begin{equation} \label{eq:3.1.2}
y=y(x_i)+f(x_i,y(x_i))(x-x_i).
\end{equation}

Setting $x=x_{i+1}=x_i+h$ in \eqref{eq:3.1.2} yields
\begin{equation} \label{eq:3.1.3}
y_{i+1}=y(x_i)+hf(x_i,y(x_i))
\end{equation}
as an approximation to $y(x_{i+1})$. Since $y(x_0)=y_0$ is known, we
can use \eqref{eq:3.1.3} with $i=0$ to compute
\[
y_1=y_0+hf(x_0,y_0).
\]
However, setting $i=1$ in \eqref{eq:3.1.3} yields
\[
y_2=y(x_1)+hf(x_1,y(x_1)),
\]
which isn't  useful, since we \emph{\color{blue}don't know} $y(x_1)$. Therefore
we replace $y(x_1)$ by its approximate value $y_1$ and redefine
\[
y_2=y_1+hf(x_1,y_1).
\]
Having computed $y_2$, we can  compute
\[
y_3=y_2+hf(x_2,y_2).
\]
In general, Euler's method starts with the known value $y(x_0)=y_0$
and computes $y_1$, $y_2$, \dots, $y_n$ successively by with the
 formula
\begin{equation} \label{eq:3.1.4}
y_{i+1}=y_i+hf(x_i,y_i),\quad 0\le i\le n-1.
\end{equation}

The next example illustrates the computational procedure
indicated in Euler's method.

\begin{example}\label{example:3.1.1} \rm
Use Euler's method with $h=0.1$ to find approximate values for the
solution of the initial value problem
\begin{equation}\label{eq:3.1.5}
y'+2y=x^3e^{-2x},\quad y(0)=1
\end{equation}
at $x=0.1,0.2,0.3$.
\end{example}

\solution
We rewrite \eqref{eq:3.1.5}  as
\[
y'=-2y+x^3e^{-2x},\quad y(0)=1,
\]
which is of the form \eqref{eq:3.1.1}, with
\[
f(x,y)=-2y+x^3e^{-2x},\ x_0=0,\mbox{\ and}\ y_0=1.
\]
Euler's method yields
\begin{eqnarray*}
y_1 & = & y_0+hf(x_0,y_0)\\
 & = & 1+(.1)f(0,1)=1+(.1)(-2)=.8,\\[4\jot]
y_2 & = & y_1+hf(x_1,y_1)\\
 & = & .8+(.1)f(.1,.8)=.8+(.1)\left(-2(.8)+(.1)^3e^{-.2}\right)=
.640081873,\\[4\jot]
y_3 & = & y_2+hf(x_2,y_2)\\
 & = & .640081873+(.1)\left(-2(.640081873)+(.2)^3e^{-.4}\right)=
.512601754.\bbox
\end{eqnarray*}

We've written the details of these computations to ensure that you
understand the procedure. However,  in the rest of the examples as
well as
 the exercises  in this chapter, we'll assume that you can use a
programmable calculator or a computer
to carry out the necessary computations.

\medskip
\boxit{Examples Illustrating The Error in Euler's Method}


\begin{example}\label{example:3.1.2} \rm
Use Euler's method with step sizes $h=0.1$, $h=0.05$, and $h=0.025$ to
find approximate values of the solution of the initial value problem
\[
y'+2y=x^3e^{-2x},\quad y(0)=1
\]
at $x=0$, $0.1$, $0.2$, $0.3$, \dots, $1.0$. Compare these approximate
values
with the values of the exact solution
\begin{equation} \label{eq:3.1.6}
y={e^{-2x}\over4}(x^4+4),
\end{equation}
which can be obtained by the method of Section~2.1. (Verify.)
\end{example}


\solution Table~\ref{table:3.1.1} shows the values of the exact solution
\eqref{eq:3.1.6} at the specified points, and the approximate values of
the solution at these points obtained by Euler's method with step
sizes $h=0.1$, $h=0.05$, and $h=0.025$. In examining this table, keep
in mind that the approximate values in the column corresponding to
$h=.05$ are actually the results of 20 steps with Euler's method. We
haven't listed the estimates of the solution obtained for
$x=0.05$, $0.15$, \dots, since there's nothing to compare them with in
the column corresponding to $h=0.1$. Similarly, the approximate values
in the column corresponding to $h=0.025$ are actually the results of
40 steps with Euler's method.

\begin{newtable} \label{table:3.1.1}
{Numerical solution of
$y'+2y=x^3e^{-2x},\ y(0)=1$, by Euler's method.}
\begin{center}
\begin{tabular}{|c|r|r|r|r|}
\hline
\multicolumn{1}{|c|}{$x$}&
\multicolumn{1}{|c|}{$h=0.1$}&
\multicolumn{1}{|c|}{$h=0.05$}&
\multicolumn{1}{|c|}{$h=0.025$}&
\multicolumn{1}{|c|}{Exact}\\ \hline
0.0 & 1.000000000 & 1.000000000 & 1.000000000 & 1.000000000 \\
0.1 & 0.800000000 & 0.810005655 & 0.814518349 & 0.818751221 \\
0.2 & 0.640081873 & 0.656266437 & 0.663635953 & 0.670588174 \\
0.3 & 0.512601754 & 0.532290981 & 0.541339495 & 0.549922980 \\
0.4 & 0.411563195 & 0.432887056 & 0.442774766 & 0.452204669 \\
0.5 & 0.332126261 & 0.353785015 & 0.363915597 & 0.373627557 \\
0.6 & 0.270299502 & 0.291404256 & 0.301359885 & 0.310952904 \\
0.7 & 0.222745397 & 0.242707257 & 0.252202935 & 0.261398947 \\
0.8 & 0.186654593 & 0.205105754 & 0.213956311 & 0.222570721 \\
0.9 & 0.159660776 & 0.176396883 & 0.184492463 & 0.192412038 \\
1.0 & 0.139778910 & 0.154715925 & 0.162003293 & 0.169169104\\
\hline
\end{tabular}
\end{center}
\end{newtable}

You can see from Table~\ref{table:3.1.1} that decreasing the step size
improves the accuracy of Euler's method. For example,
\[
y_{\mbox{\tiny exact}}(1)-y_{\mbox{\tiny approx}}(1)\approx
\left\{\begin{array}{l}
.0293\mbox{ with }h=0.1,\\
.0144\mbox{ with }h=0.05,\\
.0071\mbox{ with }h=0.025.
\end{array}\right.
\]
Based on this scanty evidence, you might guess that the error in
approximating the exact solution at a \emph{\color{blue}fixed value of} $x$ by
Euler's method is roughly halved when the step size is halved. You can
find more evidence to support this conjecture by examining
Table~\ref{table:3.1.2}, which lists the approximate values of
$y_{\mbox{\tiny exact}}-y_{\mbox{\tiny approx}}$ at
$x=0.1$, $0.2$, \dots, $1.0$.

\begin{newtable} \label{table:3.1.2}
{Errors in approximate solutions of $y'+2y=x^3e^{-2x},\
y(0)=1$, obtained by Euler's method.}
\begin{center}
\begin{tabular}{|c|r|r|r|}
\hline
\multicolumn{1}{|c|}{$x$}&
\multicolumn{1}{|c|}{$h=0.1$}&
\multicolumn{1}{|c|}{$h=0.05$}&
\multicolumn{1}{|c|}{$h=0.025$}\\ \hline
0.1 & 0.0187 & 0.0087 & 0.0042\\
0.2 & 0.0305 & 0.0143 & 0.0069\\
0.3 & 0.0373 & 0.0176 & 0.0085\\
0.4 & 0.0406 & 0.0193 & 0.0094\\
0.5 & 0.0415 & 0.0198 & 0.0097\\
0.6 & 0.0406 & 0.0195 & 0.0095\\
0.7 & 0.0386 & 0.0186 & 0.0091\\
0.8 & 0.0359 & 0.0174 & 0.0086\\
0.9 & 0.0327 & 0.0160 & 0.0079\\
1.0 & 0.0293 & 0.0144 & 0.0071\\
\hline
\end{tabular}
\end{center}
\end{newtable}

\begin{example}\label{example:3.1.3} \rm
Tables~\ref{table:3.1.3} and \ref{table:3.1.4} show analogous results
for the nonlinear initial value problem
\begin{equation} \label{eq:3.1.7}
y'=-2y^2+xy+x^2,\ y(0)=1,
\end{equation}
except in this case we can't solve \eqref{eq:3.1.7} exactly.
The results in the ``Exact'' column were obtained by using a
more accurate numerical method known as the
\href{http://www-history.mcs.st-and.ac.uk/Mathematicians/Runge.html}
{\color{blue}\itshape Runge}-\href{http://www-history.mcs.st-and.ac.uk/Mathematicians/Kutta.html}{\itshape\color{blue}Kutta\/}
method
 with a small step size. They are
exact to eight decimal places.\bbox
\end{example}


Since we think it's important in evaluating the accuracy of the
numerical methods that we'll be studying in this chapter, we
often include a column listing values of the exact solution of the
initial value problem, even if the directions in the example or
exercise don't specifically call for it. If  quotation marks are
included in the heading, the values were obtained by applying the
Runge-Kutta method in a way that's explained in Section~3.3.
If  quotation marks are not included,  the values were
obtained from a known formula for the solution.
In either case, the values are exact to eight places to the
right of the decimal point.



\begin{newtable} \label{table:3.1.3}
{Numerical solution of
$y'=-2y^2+xy+x^2,\ y(0)=1$, by Euler's method.}
\begin{center}
\begin{tabular}{|c|r|r|r|r|}
\hline
\multicolumn{1}{|c|}{$x$}&
\multicolumn{1}{|c|}{$h=0.1$}&
\multicolumn{1}{|c|}{$h=0.05$}&
\multicolumn{1}{|c|}{$h=0.025$}&
\multicolumn{1}{|c|}{``Exact''}\\ \hline
0.0 & 1.000000000 & 1.000000000 & 1.000000000 & 1.000000000 \\
0.1 & 0.800000000 & 0.821375000 & 0.829977007 & 0.837584494 \\
0.2 & 0.681000000 & 0.707795377 & 0.719226253 & 0.729641890 \\
0.3 & 0.605867800 & 0.633776590 & 0.646115227 & 0.657580377 \\
0.4 & 0.559628676 & 0.587454526 & 0.600045701 & 0.611901791 \\
0.5 & 0.535376972 & 0.562906169 & 0.575556391 & 0.587575491 \\
0.6 & 0.529820120 & 0.557143535 & 0.569824171 & 0.581942225 \\
0.7 & 0.541467455 & 0.568716935 & 0.581435423 & 0.593629526 \\
0.8 & 0.569732776 & 0.596951988 & 0.609684903 & 0.621907458 \\
0.9 & 0.614392311 & 0.641457729 & 0.654110862 & 0.666250842 \\
1.0 & 0.675192037 & 0.701764495 & 0.714151626 & 0.726015790\\
\hline
\end{tabular}
\end{center}
\end{newtable}

\begin{newtable} \label{table:3.1.4}
{Errors in approximate solutions of $y'=-2y^2+xy+x^2,\
y(0)=1$, obtained by Euler's method.}
\begin{center}
\begin{tabular}{|c|r|r|r|}
\hline
\multicolumn{1}{|c|}{$x$}&
\multicolumn{1}{|c|}{$h=0.1$}&
\multicolumn{1}{|c|}{$h=0.05$}&
\multicolumn{1}{|c|}{$h=0.025$}\\ \hline
0.1 & 0.0376 & 0.0162 &0.0076 \\
0.2 & 0.0486 & 0.0218 &0.0104 \\
0.3 & 0.0517 & 0.0238 &0.0115 \\
0.4 & 0.0523 & 0.0244 &0.0119 \\
0.5 & 0.0522 & 0.0247 &0.0121 \\
0.6 & 0.0521 & 0.0248 &0.0121 \\
0.7 & 0.0522 & 0.0249 &0.0122 \\
0.8 & 0.0522 & 0.0250 &0.0122 \\
0.9 & 0.0519 & 0.0248 &0.0121 \\
1.0 & 0.0508 & 0.0243 &0.0119 \\
\hline
\end{tabular}
\end{center}
\end{newtable}


\boxit{Truncation Error in Euler's Method}


\noindent
Consistent with the results indicated in
Tables~\ref{table:3.1.1}--\ref{table:3.1.4}, we'll now show that under
reasonable assumptions on $f$ there's a constant $K$ such that
the error in approximating the solution of the initial value problem
\[
y'=f(x,y),\quad y(x_0)=y_0,
\]
at a given point $b>x_0$ by Euler's method with step size
$h=(b-x_0)/n$ satisfies the inequality
\[
|y(b)-y_n|\le Kh,
\]
where $K$ is a constant independent of $n$.


There are two sources of error (not counting roundoff) in Euler's
method:
\begin{numberlist}
\item % (1)
The error committed in approximating the integral curve by the tangent
line \eqref{eq:3.1.2} over the interval $[x_i,x_{i+1}]$.
\item % (2)
The error committed in replacing $y(x_i)$ by $y_i$ in
\eqref{eq:3.1.2} and using \eqref{eq:3.1.4} rather than \eqref{eq:3.1.2} to
compute $y_{i+1}$.
 \end{numberlist}

Euler's method assumes that $y_{i+1}$ defined in \eqref{eq:3.1.2}  is
an approximation to $y(x_{i+1})$. We call the error in this
approximation the \emph{\color{blue}local truncation error at the $i$th step},
and denote it by $T_i$;   thus,
\begin{equation} \label{eq:3.1.8}
T_i=y(x_{i+1})-y(x_i)-hf(x_i,y(x_i)).
\end{equation}
We'll now use
\href{http://www-history.mcs.st-and.ac.uk/Mathematicians/Taylor.html}
{\color{blue}\itshape Taylor's theorem\/}
 to estimate $T_i$, assuming for
simplicity that $f$, $f_x$, and $f_y$ are continuous and bounded for
all $(x,y)$. Then $y''$ exists and is bounded on $[x_0,b]$. To see
this, we differentiate
\[
y'(x)=f(x,y(x))
\]
to obtain
\begin{eqnarray*}
y''(x) & = & f_x(x,y(x))+f_y(x,y(x))y'(x)\\
 & = & f_x(x,y(x))+f_y(x,y(x))f(x,y(x)).
\end{eqnarray*}
Since we assumed that $f$, $f_x$ and $f_y$ are bounded, there's
a constant $M$ such that
\[
|f_x(x,y(x))+f_y(x,y(x))f(x,y(x))|\le M,\quad x_0<x<b,
\]
which implies that
\begin{equation} \label{eq:3.1.9}
|y''(x)|\le M,\quad x_0<x<b.
\end{equation}
Since $x_{i+1}=x_i+h$, Taylor's theorem implies that
\[
y(x_{i+1})=y(x_i)+hy'(x_i)+{h^2\over2}y''(\tilde x_i),
\]
where $\tilde x_i$ is some number between $x_i$ and $x_{i+1}$.
Since $y'(x_i)=f(x_i,y(x_i))$  this can be written as
\[
y(x_{i+1})=y(x_i)+hf(x_i,y(x_i))+{h^2\over2}y''(\tilde x_i),
\]
or, equivalently,
\[
y(x_{i+1})-y(x_i)-hf(x_i,y(x_i))={h^2\over2}y''(\tilde x_i).
\]
Comparing this with \eqref{eq:3.1.8} shows that
\[
T_i={h^2\over2}y''(\tilde x_i).
\]
Recalling \eqref{eq:3.1.9}, we can establish the bound
\begin{equation} \label{eq:3.1.10}
|T_i|\le{Mh^2\over2},\quad 1\le i\le n.
\end{equation}
Although it may be difficult to determine the constant $M$, what is
important is that there's an $M$ such that \eqref{eq:3.1.10} holds.
 We say that the
local truncation error of Euler's method is \emph{\color{blue}of order} $h^2$,
which we write as $O(h^2)$.

Note that the magnitude of the local truncation error in
Euler's method is determined by the second derivative $y''$ of the
solution of the initial value problem. Therefore
 the local truncation error will be larger where $|y''|$ is large,
or smaller where $|y''|$ is small.

Since the local truncation error for Euler's method is $O(h^2)$, it's
reasonable to expect that halving $h$ reduces the local truncation
error by a factor of 4. This is true, but
 halving the step size also requires twice as many steps to
approximate the solution at a given point. To analyze the overall
effect of truncation error in Euler's method, it's useful to  derive
an equation relating the errors
\[
e_{i+1}=y(x_{i+1})-y_{i+1}\mbox{\quad and \quad} e_i=y(x_i)-y_i.
\]
To this end, recall that
\begin{equation} \label{eq:3.1.11}
y(x_{i+1})=y(x_i)+hf(x_i,y(x_i))+T_i
\end{equation}
and
\begin{equation} \label{eq:3.1.12}
y_{i+1}=y_i+hf(x_i,y_i).
\end{equation}
Subtracting \eqref{eq:3.1.12} from \eqref{eq:3.1.11} yields
\begin{equation} \label{eq:3.1.13}
e_{i+1}=e_i+h\left[f(x_i,y(x_i))-f(x_i,y_i)\right]+T_i.
\end{equation}
The last term on the right is the local truncation error at the $i$th
step. The other terms reflect the way errors made at \emph{\color{blue}previous
steps} affect $e_{i+1}$. Since $|T_i|\le Mh^2/2$, we see from
\eqref{eq:3.1.13} that
\begin{equation} \label{eq:3.1.14}
|e_{i+1}|\le |e_i|+h|f(x_i,y(x_i))-f(x_i,y_i)|+{Mh^2\over2}.
\end{equation}
Since we assumed that $f_y$ is continuous and bounded, the mean
value theorem implies that
\[
f(x_i,y(x_i))-f(x_i,y_i)=f_y(x_i,y_i^*)(y(x_i)-y_i)=f_y(x_i,y_i^*)e_i,
\]
where $y_i^*$ is between $y_i$ and $y(x_i)$. Therefore
\[
|f(x_i,y(x_i))-f(x_i,y_i)|\le R|e_i|
\]
for some constant $R$. From this and \eqref{eq:3.1.14},
\begin{equation} \label{eq:3.1.15}
|e_{i+1}|\le (1+Rh)|e_i|+{Mh^2\over2},\quad 0\le i\le n-1.
\end{equation}
For convenience, let $C=1+Rh$. Since $e_0=y(x_0)-y_0=0$, applying
\eqref{eq:3.1.15} repeatedly yields
\begin{eqnarray}
|e_1| & \le & {Mh^2\over2}\nonumber\\
|e_2| & \le & C|e_1|+{Mh^2\over2}\le(1+C){Mh^2\over2}\nonumber\\
|e_3| & \le & C|e_2|+{Mh^2\over2}\le(1+C+C^2){Mh^2\over2}\nonumber\\
 & \vdots &\nonumber \\|e_n| & \le &
C|e_{n-1}|+{Mh^2\over2}\le(1+C+\cdots+C^{n-1}){Mh^2\over2}.\label{eq:3.1.16}
\end{eqnarray}

Recalling the formula for the sum of a geometric series, we see that
\[
1+C+\cdots+C^{n-1}={1-C^n\over 1-C}={(1+Rh)^n-1\over Rh}
\]
(since $C=1+Rh$). From this and \eqref{eq:3.1.16},
\begin{equation} \label{eq:3.1.17}
|y(b)-y_n|=|e_n|\le{(1+Rh)^n-1\over R}{Mh\over2}.
\end{equation}
Since Taylor's theorem implies that
\[
1+Rh<e^{Rh}
\]
(verify),
\[
(1+Rh)^n<e^{nRh}=e^{R(b-x_0)}\quad (\mbox{since }nh=b-x_0).
\]
This and \eqref{eq:3.1.17} imply that
\begin{equation} \label{eq:3.1.18}
|y(b)-y_n|\le Kh,
\end{equation}
with
\[
K=M{e^{R(b-x_0)}-1\over2R}.
\]
Because of \eqref{eq:3.1.18} we say that the
\href{http://www-history.mcs.st-and.ac.uk/Mathematicians/Euler.html}
\emph{\color{blue}global truncation
error of Euler's method is of order} $h$, which we write as $O(h)$.

\boxit{Semilinear Equations and Variation of Parameters}


\noindent
An equation that can be written in the form
\begin{equation} \label{eq:3.1.19}
y'+p(x)y=h(x,y)
\end{equation}
with $p\not\equiv0$ is said to be \emph{\color{blue}semilinear}. (Of course,
\eqref{eq:3.1.19} is linear if $h$ is independent of $y$.) One way to
apply Euler's method to an initial value problem
\begin{equation} \label{eq:3.1.20}
y'+p(x)y=h(x,y),\quad y(x_0)=y_0
\end{equation}
for \eqref{eq:3.1.19} is to think of it  as
\[
y'=f(x,y),\quad y(x_0)=y_0,
\]
where
\[
f(x,y)=-p(x)y+h(x,y).
\]
However, we can also start by applying variation of parameters to
\eqref{eq:3.1.20}, as in Sections~2.1\ and 2.4;   thus, we
write the solution of \eqref{eq:3.1.20} as $y=uy_1$, where $y_1$ is a
nontrivial solution of the complementary equation $y'+p(x)y=0$. Then
$y=uy_1$ is a solution of \eqref{eq:3.1.20} if and only if $u$ is a
solution of the initial value problem
\begin{equation} \label{eq:3.1.21}
u'=h(x,uy_1(x))/y_1(x),\quad u(x_0)=y(x_0)/y_1(x_0).
\end{equation}
We can apply Euler's method to obtain approximate values
$u_0$, $u_1$, \dots, $u_n$ of this initial value problem, and then take
\[
y_i=u_iy_1(x_i)
\]
as approximate values of the solution of \eqref{eq:3.1.20}.
We'll call this procedure the
\href{http://www-history.mcs.st-and.ac.uk/Mathematicians/Euler.html}
{\color{blue}\itshape  Euler semilinear method\/}.

The next two examples show that the Euler and Euler semilinear
methods may yield drastically different results.

\enlargethispage{3in}

\begin{example}\label{example:3.1.4} \rm
In Example~\ref{example:2.1.7} we had to leave the solution of the initial
value
problem
\begin{equation} \label{eq:3.1.22}
y'-2xy=1,\quad y(0)=3
\end{equation}
in the  form
\begin{equation} \label{eq:3.1.23}
y=e^{x^2}\left(3 +\int^x_0 e^{-t^2}dt\right)
\end{equation}
\newpage
\noindent
because it was impossible to evaluate this integral exactly in terms
of elementary functions.
Use step sizes $h=0.2$, $h=0.1$, and $h=0.05$ to find approximate
values
of the solution of \eqref{eq:3.1.22} at $x=0$, $0.2$, $0.4$, $0.6$,
\dots, $2.0$ by
{\bf (a)} Euler's method;   {\bf (b)} the Euler semilinear method.
\end{example}


\solutionpart{a}
Rewriting \eqref{eq:3.1.22} as
\begin{equation} \label{eq:3.1.24}
y'=1+2xy,\quad y(0)=3
\end{equation}
and applying Euler's method with $f(x,y)=1+2xy$ yields the results
shown in Table~\ref{table:3.1.5}. Because of the large differences
between the estimates obtained for the three values of $h$, it
would be clear that these results are useless even if the ``exact''
values were not included in the table.

\begin{newtable} \label{table:3.1.5}
{Numerical solution of
$y'-2xy=1,\ y(0)=3$, with Euler's method.}
\begin{center}
\begin{tabular}{|c|r|r|r|r|}
\hline
\multicolumn{1}{|c|}{$x$}&
\multicolumn{1}{|c|}{$h=0.2$}&
\multicolumn{1}{|c|}{$h=0.1$}&
\multicolumn{1}{|c|}{$h=0.05$}&
\multicolumn{1}{|c|}{``Exact''}\\ \hline
0.0 &  3.000000000 & 3.000000000 & 3.000000000     &   3.000000000 \\
0.2 &  3.200000000 & 3.262000000 &3.294348537      &   3.327851973 \\
0.4 &  3.656000000 & 3.802028800 & 3.881421103     &   3.966059348 \\
0.6 &  4.440960000 & 4.726810214 & 4.888870783     &   5.067039535 \\
0.8 &  5.706790400 & 6.249191282 & 6.570796235     &   6.936700945 \\
1.0 &  7.732963328 & 8.771893026 & 9.419105620     &  10.184923955 \\
1.2 & 11.026148659 &13.064051391 &14.405772067     &  16.067111677 \\
1.4 & 16.518700016 &20.637273893 & 23.522935872    &  27.289392347 \\
1.6 & 25.969172024 &34.570423758 & 41.033441257    &  50.000377775 \\
1.8 & 42.789442120 &61.382165543 & 76.491018246    &  98.982969504 \\
2.0 & 73.797840446 & 115.440048291 & 152.363866569 & 211.954462214 \\
\hline
\end{tabular}
\end{center}
\end{newtable}


It's easy to see why Euler's method yields such poor results.
Recall that the constant $M$ in
\eqref{eq:3.1.10} -- which plays an important role in determining the
local truncation error in Euler's method -- must be an upper bound for
the values of the second derivative $y''$ of the solution of the
initial
value problem \eqref{eq:3.1.22} on $(0,2)$. The problem is that $y''$
assumes very large values on this interval. To see this, we
differentiate \eqref{eq:3.1.24} to obtain
\[
y''(x)=2y(x)+2xy'(x)=2y(x)+2x(1+2xy(x))=2(1+2x^2)y(x)+2x,
\]
where the second equality follows again from \eqref{eq:3.1.24}.
Since \eqref{eq:3.1.23} implies that $y(x)>3e^{x^2}$ if $x>0$,
\[
y''(x)>6(1+2x^2)e^{x^2}+2x,\quad x>0.
\]
For example, letting $x=2$ shows that $y''(2)>2952$.

\solutionpart{b}
Since $y_1=e^{x^2}$ is a solution of the complementary equation
$y'-2xy=0$, we can  apply the Euler semilinear method to
\eqref{eq:3.1.22}, with
\[
y=ue^{x^2}\mbox{\quad and \quad}
u'=e^{-x^2},\quad u(0)=3.
\]
The results listed in Table~\ref{table:3.1.6} are clearly better than
those obtained by Euler's method.

\begin{newtable} \label{table:3.1.6}
{Numerical solution of
$y'-2xy=1,\ y(0)=3$, by the Euler semilinear method.}
\begin{center}
\begin{tabular}{|c|r|r|r|r|}
\hline
\multicolumn{1}{|c|}{$x$}&
\multicolumn{1}{|c|}{$h=0.2$}&
\multicolumn{1}{|c|}{$h=0.1$}&
\multicolumn{1}{|c|}{$h=0.05$}&
\multicolumn{1}{|c|}{``Exact''}\\ \hline
0.0 &  3.000000000  &   3.000000000 &   3.000000000 &   3.000000000 \\
0.2 &  3.330594477  &   3.329558853 &   3.328788889 &   3.327851973 \\
0.4 &  3.980734157  &   3.974067628 &   3.970230415 &   3.966059348 \\
0.6 &  5.106360231  &   5.087705244 &   5.077622723 &   5.067039535 \\
0.8 &  7.021003417  &   6.980190891 &   6.958779586 &   6.936700945 \\
1.0 & 10.350076600  &  10.269170824 &  10.227464299 &  10.184923955 \\
1.2 & 16.381180092  &  16.226146390 &  16.147129067 &  16.067111677 \\
1.4 & 27.890003380  &  27.592026085 &  27.441292235 &  27.289392347 \\
1.6 & 51.183323262  &  50.594503863 &  50.298106659 &  50.000377775 \\
1.8 &101.424397595  & 100.206659076 &  99.595562766 &  98.982969504 \\
2.0 &217.301032800  & 214.631041938 & 213.293582978 & 211.954462214\\
\hline
\end{tabular}
\end{center}
\end{newtable}


We can't give a general procedure for determining in advance
whether Euler's method or the semilinear Euler method will produce
better results for a given semilinear initial value problem
\eqref{eq:3.1.19}.
 As a rule of thumb, the Euler semilinear method will
yield better results than Euler's method if $|u''|$ is small on
$[x_0,b]$, while Euler's method  yields better results if $|u''|$
is large on $[x_0,b]$.
 In many cases the results obtained by the two methods don't
 differ appreciably. However, we propose the an intuitive
way to decide which is the better method: Try both methods with
multiple step sizes, as we did in Example~\ref{example:3.1.4},
 and accept the results obtained by the method for
which the approximations change less as the step size decreases.

\begin{example}\label{example:3.1.5} \rm
Applying Euler's method with step sizes $h=0.1$, $h=0.05$, and
$h=0.025$ to the initial value problem
\begin{equation}\label{eq:3.1.25}
y'-2y={x\over1+y^2},\quad y(1)=7
\end{equation}
on $[1,2]$ yields the results in
Table~\ref{table:3.1.7}. Applying the Euler semilinear method
with
\[
y=ue^{2x}\mbox{\quad and \quad}u'={xe^{-2x}\over1+u^2e^{4x}},\quad
u(1)=7e^{-2}
\]
 yields the results in Table~\ref{table:3.1.8}.
Since the latter are clearly less dependent on step size
than the former, we conclude that the Euler semilinear method
is better than Euler's method for \eqref{eq:3.1.25}. This conclusion
is supported by comparing the approximate results obtained by the
two methods with the ``exact'' values of the solution.
\end{example}


\begin{newtable} \label{table:3.1.7}
{Numerical solution of
$y'-2y=x/(1+y^2),\ y(1)=7$, by Euler's method.}
\begin{center}
\begin{tabular}{|c|r|r|r|r|}
\hline
\multicolumn{1}{|c|}{$x$}&
\multicolumn{1}{|c|}{$h=0.1$}&
\multicolumn{1}{|c|}{$h=0.05$}&
\multicolumn{1}{|c|}{$h=0.025$}&
\multicolumn{1}{|c|}{``Exact''} \\ \hline
1.0 &  7.000000000 &  7.000000000 &  7.000000000 &  7.000000000\\
1.1 &  8.402000000 &  8.471970569 &  8.510493955 &  8.551744786\\
1.2 & 10.083936450 & 10.252570169 & 10.346014101 & 10.446546230\\
1.3 & 12.101892354 & 12.406719381 & 12.576720827 & 12.760480158\\
1.4 & 14.523152445 & 15.012952416 & 15.287872104 & 15.586440425\\
1.5 & 17.428443554 & 18.166277405 & 18.583079406 & 19.037865752\\
1.6 & 20.914624471 & 21.981638487 & 22.588266217 & 23.253292359\\
1.7 & 25.097914310 & 26.598105180 & 27.456479695 & 28.401914416\\
1.8 & 30.117766627 & 32.183941340 & 33.373738944 & 34.690375086\\
1.9 & 36.141518172 & 38.942738252 & 40.566143158 & 42.371060528\\
2.0 & 43.369967155 & 47.120835251 & 49.308511126 & 51.752229656\\
\hline
\end{tabular}
\end{center}
\end{newtable}


\begin{newtable} \label{table:3.1.8}
{Numerical solution of
$y'-2y=x/(1+y^2),\ y(1)=7$, by the Euler semilinear method.}
\begin{center}
\begin{tabular}{|c|r|r|r|r|r|}
\hline
\multicolumn{1}{|c|}{$x$}&
\multicolumn{1}{|c|}{$h=0.1$}&
\multicolumn{1}{|c|}{$h=0.05$}&
\multicolumn{1}{|c|}{$h=0.025$}&
\multicolumn{1}{|c|}{``Exact''} \\ \hline
1.0 &  7.000000000 &  7.000000000 &  7.000000000 &  7.000000000\\
1.1 &  8.552262113 &  8.551993978 &  8.551867007 &  8.551744786\\
1.2 & 10.447568674 & 10.447038547 & 10.446787646 & 10.446546230\\
1.3 & 12.762019799 & 12.761221313 & 12.760843543 & 12.760480158\\
1.4 & 15.588535141 & 15.587448600 & 15.586934680 & 15.586440425\\
1.5 & 19.040580614 & 19.039172241 & 19.038506211 & 19.037865752\\
1.6 & 23.256721636 & 23.254942517 & 23.254101253 & 23.253292359\\
1.7 & 28.406184597 & 28.403969107 & 28.402921581 & 28.401914416\\
1.8 & 34.695649222 & 34.692912768 & 34.691618979 & 34.690375086\\
1.9 & 42.377544138 & 42.374180090 & 42.372589624 & 42.371060528\\
2.0 & 51.760178446 & 51.756054133 & 51.754104262 & 51.752229656\\
\hline
\end{tabular}
\end{center}
\end{newtable}

\begin{example}{\color{blue}}\label{example:3.1.6} \rm
Applying Euler's method with step sizes $h=0.1$, $h=0.05$, and
$h=0.025$ to  the initial value problem
\begin{equation}\label{eq:3.1.26}
y'+3x^2y=1+y^2,\quad y(2)=2
\end{equation}
on  $[2,3]$ yields the results in
Table~\ref{table:3.1.9}. Applying the Euler semilinear method with
\[
y=ue^{-x^3}\mbox{\quad and \quad}u'=e^{x^3}(1+u^2e^{-2x^3}),\quad
u(2)=2e^8
\]
yields the results in Table~\ref{table:3.1.10}. Noting the close
agreement among the  three columns of Table~\ref{table:3.1.9} (at
least for larger values of $x$) and the lack of any such agreement
among the columns of Table~\ref{table:3.1.10}, we conclude that
Euler's method is better than the Euler semilinear method for
\eqref{eq:3.1.26}. Comparing the results with the exact values
supports this conclusion.
\end{example}


\begin{newtable} \label{table:3.1.9}
{Numerical solution of
$y'+3x^2y=1+y^2,\quad y(2)=2$, by Euler's method.}
\begin{center}
\begin{tabular}{|c|r|r|r|r|r|}
\hline
\multicolumn{1}{|c|}{$x$}&
\multicolumn{1}{|c|}{$h=0.1$}&
\multicolumn{1}{|c|}{$h=0.05$}&
\multicolumn{1}{|c|}{$h=0.025$}&
\multicolumn{1}{|c|}{``Exact''}\\ \hline
2.0 &2.000000000 & 2.000000000 & 2.000000000 & 2.000000000 \\
2.1 &0.100000000 & 0.493231250 & 0.609611171 & 0.701162906 \\
2.2 &0.068700000 & 0.122879586 & 0.180113445 & 0.236986800 \\
2.3 &0.069419569 & 0.070670890 & 0.083934459 & 0.103815729 \\
2.4 &0.059732621 & 0.061338956 & 0.063337561 & 0.068390786 \\
2.5 &0.056871451 & 0.056002363 & 0.056249670 & 0.057281091 \\
2.6 &0.050560917 & 0.051465256 & 0.051517501 & 0.051711676 \\
2.7 &0.048279018 & 0.047484716 & 0.047514202 & 0.047564141 \\
2.8 &0.042925892 & 0.043967002 & 0.043989239 & 0.044014438 \\
2.9 &0.042148458 & 0.040839683 & 0.040857109 & 0.040875333 \\
3.0 &0.035985548 & 0.038044692 & 0.038058536 & 0.038072838 \\
\hline
\end{tabular}
\end{center}
\end{newtable}

\begin{newtable} \label{table:3.1.10}
{Numerical solution of
$y'+3x^2y=1+y^2,\quad y(2)=2$, by the Euler semilinear method.}
\begin{center}
\begin{tabular}{|c|r|r|r|r|r|}
\hline
\multicolumn{1}{|c|}{$x$}&
\multicolumn{1}{|c|}{$h=0.1$}&
\multicolumn{1}{|c|}{$h=0.05$}&
\multicolumn{1}{|c|}{$h=0.025$}&
\multicolumn{1}{|c|}{``Exact''}\\ \hline
$x$ & $h=0.1$ & $h=0.05$ & $h=0.025$ &$h=.0125$\\ \hline
2.0 &2.000000000 & 2.000000000 & 2.000000000 & 2.000000000 \\
2.1 &0.708426286 & 0.702568171 & 0.701214274 & 0.701162906 \\
2.2 &0.214501852 & 0.222599468 & 0.228942240 & 0.236986800 \\
2.3 &0.069861436 & 0.083620494 & 0.092852806 & 0.103815729 \\
2.4 &0.032487396 & 0.047079261 & 0.056825805 & 0.068390786 \\
2.5 &0.021895559 & 0.036030018 & 0.045683801 & 0.057281091 \\
2.6 &0.017332058 & 0.030750181 & 0.040189920 & 0.051711676 \\
2.7 &0.014271492 & 0.026931911 & 0.036134674 & 0.047564141 \\
2.8 &0.011819555 & 0.023720670 & 0.032679767 & 0.044014438 \\
2.9 &0.009776792 & 0.020925522 & 0.029636506 & 0.040875333 \\
3.0 &0.008065020 & 0.018472302 & 0.026931099 & 0.038072838 \\
\hline
\end{tabular}
\end{center}
\end{newtable}



In the next two sections we'll study other numerical methods for
solving initial value problems,  called the
\href{http://www-history.mcs.st-and.ac.uk/Mathematicians/Euler.html}
 {\color{blue}\itshape improved Euler method}, the \emph{\color{blue}midpoint method},
\href{http://www-history.mcs.st-and.ac.uk/Mathematicians/Heun.html}
{\color{blue}\itshape Heun's method}
and the
\href{http://www-history.mcs.st-and.ac.uk/Mathematicians/Runge.html}
{\color{blue}\itshape
Runge}-\href{http://www-history.mcs.st-and.ac.uk/Mathematicians/Kutta.html}
{\itshape\color{blue}Kutta method\/}.
 If the initial value problem is semilinear as
in \eqref{eq:3.1.19},  we also have the option of using variation of
parameters and then applying the given numerical method to the initial
value problem \eqref{eq:3.1.21} for $u$. By analogy with the terminology
used here, we'll call the resulting procedure
\href{http://www-history.mcs.st-and.ac.uk/Mathematicians/Euler.html}
{\color{blue}\itshape the improved
Euler semilinear method}, the \emph{\color{blue}midpoint semilinear
method},
\href{http://www-history.mcs.st-and.ac.uk/Mathematicians/Heun.html}
{\itshape\color{blue}Heun's
semilinear method} or the
\href{http://www-history.mcs.st-and.ac.uk/Mathematicians/Runge.html}
{\color{blue}\itshape Runge}-\href{http://www-history.mcs.st-and.ac.uk/Mathematicians/Kutta.html}{\itshape\color{blue}Kutta semilinear method}, as
the case may be.

\exercises
 You may want to save the results of these
exercises, since we'll revisit in the next two
sections.
In Exercises~\ref{exer:3.1.1}--\ref{exer:3.1.5} use Euler's method to find
approximate values of the solution of the given initial value problem
at the points $x_i=x_0+ih$, where $x_0$ is the point wher the
initial condition is imposed and $i=1$, $2$, $3$. The purpose of these
exercises is to familiarize you with the computational procedure of
Euler's method.

\begin{exerciselist}
\item\label{exer:3.1.1} \Cex
$y'=2x^2+3y^2-2,\quad y(2)=1;\quad h=0.05$
\item\label{exer:3.1.2} \Cex
$y'=y+\sqrt{x^2+y^2},\quad y(0)=1;\quad h=0.1$
\item\label{exer:3.1.3} \Cex
$y'+3y=x^2-3xy+y^2,\quad y(0)=2;\quad h=0.05$
\item\label{exer:3.1.4} \Cex
$y'=\dst{1+x\over1-y^2},\quad y(2)=3;\quad h=0.1$
\item\label{exer:3.1.5} \Cex
$y'+x^2y=\sin xy,\quad y(1)=\pi;\quad h=0.2$

\item\label{exer:3.1.6} \Cex
Use Euler's method with step sizes $h=0.1$, $h=0.05$, and $h=0.025$ to
find approximate values of the solution of the initial value problem
\[
y'+3y=7e^{4x},\quad y(0)=2
\]
at $x=0$, $0.1$, $0.2$, $0.3$, \dots, $1.0$. Compare these approximate
values with
the values of the exact solution $y=e^{4x}+e^{-3x}$, which can be
obtained by the method of Section~2.1. Present your results
in a table like Table~\ref{table:3.1.1}.


\item\label{exer:3.1.7} \Cex
Use Euler's method with step sizes $h=0.1$, $h=0.05$, and $h=0.025$ to
find approximate values of the solution of the initial value problem
\[
y'+{2\over x}y={3\over x^3}+1,\quad y(1)=1
\]
at $x=1.0$, $1.1$, $1.2$, $1.3$, \dots, $2.0$. Compare these approximate
values
with the values of the exact solution
\[
y={1\over3x^2}(9\ln x+x^3+2),
\]
which can be obtained by the method of Section~2.1. Present
your results in a table like Table~\ref{table:3.1.1}.

\item\label{exer:3.1.8} \Cex
Use Euler's method with step sizes $h=0.05$, $h=0.025$, and $h=0.0125$ to
find approximate values of the solution of the initial value problem
\[
y'={y^2+xy-x^2\over x^2},\quad y(1)=2
\]
at $x=1.0$, $1.05$, $1.10$, $1.15$, \dots, $1.5$. Compare these approximate
values
with the values of the exact solution
\[
y={x(1+x^2/3)\over1-x^2/3}
\]
obtained in Example~\ref{example:2.4.3}. Present your results in a table
like Table~\ref{table:3.1.1}.


\item\label{exer:3.1.9} \Cex
In Example~\ref{example:2.2.3} it was shown that
\[
y^5+y=x^2+x-4
\]
is an implicit solution of  the initial value problem
\[
y'={2x+1\over5y^4+1},\quad y(2)=1.
\eqno{\rm(A)}
\]
Use Euler's method with step sizes $h=0.1$, $h=0.05$, and $h=0.025$ to
find approximate values of the solution of (A) at
$x=2.0$, $2.1$, $2.2$, $2.3$, \dots, $3.0$. Present your results in tabular
form.
To
check the error in these approximate values, construct another table
of values of the residual
\[
  R(x,y)=y^5+y-x^2-x+4
\]
for each value of $(x,y)$ appearing in the first table.


\item\label{exer:3.1.10} \Cex
You can see from Example~\ref{example:2.5.1} that
\[
x^4y^3+x^2y^5+2xy=4
\]
is an implicit solution of  the initial value problem
\[
y'=-{4x^3y^3+2xy^5+2y\over3x^4y^2+5x^2y^4+2x},\quad y(1)=1.
\eqno{\rm(A)}
\]
Use Euler's method with step sizes $h=0.1$, $h=0.05$, and $h=0.025$ to
find approximate values of the solution of (A) at
$x=1.0$, $1.1$, $1.2$, $1.3$, \dots, $2.0$. Present your results in tabular
form.
To check the error in these approximate values, construct another
table of values of the residual
\[
R(x,y)=x^4y^3+x^2y^5+2xy-4
\]
for each value of $(x,y)$ appearing in the first table.

\item\label{exer:3.1.11} \Cex
Use Euler's method with step sizes $h=0.1$, $h=0.05$, and $h=0.025$ to
find approximate values of the solution of the initial value problem
\[
(3y^2+4y)y'+2x+\cos x=0, \quad y(0)=1; \mbox{\;
(Exercise~2.2.~\hspace*{-3pt}\ref{exer:2.2.13})}
\]
at $x=0$, $0.1$, $0.2$, $0.3$, \dots, $1.0$.


\item\label{exer:3.1.12} \Cex
Use Euler's method with step sizes $h=0.1$, $h=0.05$, and $h=0.025$ to
find approximate values of the solution of the initial value problem
\[
y'+{(y+1)(y-1)(y-2)\over x+1}=0, \quad y(1)=0
\mbox{\; (Exercise~2.2.~\hspace*{-3pt}\ref{exer:2.2.14})}
\]
at $x=1.0$, $1.1$, $1.2$, $1.3$, \dots, $2.0$.


\item\label{exer:3.1.13} \Cex
Use Euler's method and the Euler semilinear method with step sizes
$h=0.1$, $h=0.05$, and $h=0.025$ to find approximate values of the
solution of the initial value problem
\[
y'+3y=7e^{-3x},\quad y(0)=6
\]
at $x=0$, $0.1$, $0.2$, $0.3$, \dots, $1.0$. Compare these approximate
values with
the values of the exact solution $y=e^{-3x}(7x+6)$, which can be
obtained by the method of Section~2.1. Do you notice anything
special about the results? Explain.

\exercisetext{
The linear initial value problems in
Exercises~\ref{exer:3.1.14}--\ref{exer:3.1.19} can't be solved exactly in
terms of known elementary functions. In each exercise, use
Euler's method and the Euler semilinear methods
with the indicated step sizes to find approximate values of the
solution of the given initial value problem at 11 equally spaced
points (including the endpoints) in the interval.}

\item\label{exer:3.1.14} \Cex
$y'-2y=\dst{1\over1+x^2},\quad y(2)=2$;  \;
$h=0.1,0.05,0.025$ on $[2,3]$



\item\label{exer:3.1.15} \Cex
$y'+2xy=x^2,\quad y(0)=3$\;
(Exercise~2.1.~\hspace*{-3pt}\ref{exer:2.1.38});  \;
$h=0.2,0.1,0.05$ on $[0,2]$


\item\label{exer:3.1.16} \Cex
$\dst{y'+{1\over x}y={\sin x\over x^2},\quad y(1)=2}$;\;
(Exercise~2.1.~\hspace*{-3pt}\ref{exer:2.1.39});  \;
$h=0.2,0.1,0.05$ on $[1,3]$


\item\label{exer:3.1.17} \Cex
$\dst{y'+y={e^{-x}\tan x\over x},\quad y(1)=0}$;\;
(Exercise~2.1.~\hspace*{-3pt}\ref{exer:2.1.40});  \;
$h=0.05,0.025,0.0125$ on $[1,1.5]$



\item\label{exer:3.1.18} \Cex
$\dst{y'+{2x\over 1+x^2}y={e^x\over (1+x^2)^2}, \quad y(0)=1}$;
\; (Exercise~2.1.~\hspace*{-3pt}~\hspace*{-3pt}\ref{exer:2.1.41});  \;
$h=0.2,0.1,0.05$ on $[0,2]$


\item\label{exer:3.1.19} \Cex
$xy'+(x+1)y=e^{x^2},\quad y(1)=2$;\;
(Exercise~2.1.~\hspace*{-3pt}\ref{exer:2.1.42});  \;
$h=0.05,0.025,0.0125$ on $[1,1.5]$


\exercisetext{
In Exercises~\ref{exer:3.1.20}--\ref{exer:3.1.22},
use Euler's method and
the Euler semilinear method
with the indicated  step sizes to find approximate values
of the solution of the given initial value problem at 11 equally
spaced points (including the endpoints) in the interval.}

\item\label{exer:3.1.20} \Cex
$y'+3y=xy^2(y+1),\quad y(0)=1$;  \;
$h=0.1,0.05,0.025$ on $[0,1]$

\item\label{exer:3.1.21} \Cex
$\dst{y'-4y={x\over y^2(y+1)},\quad y(0)=1}$;  \;
 $h=0.1,0.05,0.025$ on $[0,1]$

\item\label{exer:3.1.22} \Cex
$\dst{y'+2y={x^2\over1+y^2},\quad y(2)=1}$;  \;
 $h=0.1,0.05,0.025$ on $[2,3]$

\item\label{exer:3.1.23}
{\sc Numerical Quadrature.}
The fundamental theorem of  calculus says that if $f$
is continuous on a closed interval $[a,b]$ then it has an
antiderivative $F$ such that $F'(x)=f(x)$ on $[a,b]$  and
\[
\int_a^bf(x)\,dx=F(b)-F(a).
\eqno{\rm(A)}
\]
This solves the problem of evaluating a definite integral if
the integrand $f$ has an antiderivative that can be found and
evaluated easily. However, if $f$
doesn't have this property,  (A) doesn't  provide a useful
way to evaluate the definite integral. In this case we must resort to
approximate methods. There's a class of such methods called
\emph{\color{blue}numerical quadrature}, where the approximation takes the
form
\[
\int_a^bf(x)\,dx\approx \sum_{i=0}^n c_if(x_i),
\eqno{\rm(B)}
\]
where $a=x_0<x_1<\cdots<x_n=b$ are suitably chosen points and
$c_0$, $c_1$, \dots, $c_n$
are suitably chosen constants. We call (B)  a \emph{\color{blue}quadrature
formula}.
\begin{alist}
\item % (a)
Derive the quadrature formula
\[
\int_a^bf(x)\,dx\approx
h\sum_{i=0}^{n-1}f(a+ih)\mbox{\quad(where $h=(b-a)/n)$}
\eqno{\rm(C)}
\]
by applying Euler's method to the initial value problem
\[
y'=f(x),\quad y(a)=0.
\]
\item % (b)
The quadrature formula (C) is sometimes called \emph{\color{blue}the left
rectangle rule}. Draw a figure that justifies this terminology.
\item % (c)
\Lex
For several choices of $a$, $b$, and $A$,
 apply (C) to $f(x)=A$ with $n = 10,20,40,80,160,320$.
Compare your results with the exact answers
and explain what you find.
\item % (d)
\Lex
For several choices of $a$, $b$,  $A$, and $B$,
 apply (C) to $f(x)=A+Bx$ with $n=10$, $20$, $40$, $80$, $160$, $320$.
 Compare your results with the exact answers
and explain what you find.
\end{alist}
\end{exerciselist}





\setcounter{table}{0}
\section{The Improved Euler Method and Related Methods}

 In Section~3.1 we saw that the global truncation error of
Euler's method is $O(h)$, which would seem to imply that we can
achieve arbitrarily accurate results with Euler's method by simply
choosing the step size sufficiently small. However, this
isn't  a good idea, for two reasons. First,
after a certain point decreasing the step size will
increase roundoff errors to the point where the accuracy
will deteriorate rather than improve.
 The second
and more important reason is that in most applications of numerical
methods to an initial value problem
\begin{equation} \label{eq:3.2.1}
y'=f(x,y),\quad y(x_0)=y_0,
\end{equation}
the expensive part of the computation is the evaluation of $f$.
Therefore we want methods that give good results for a given number
of such evaluations. This is what motivates us to look for numerical
methods better than Euler's.

To clarify this point, suppose  we want to approximate the value of $e$
by applying Euler's method to the initial value problem
\[
y'=y,\quad y(0)=1,\quad\mbox{(with solution $y=e^x$)}
\]
on $[0,1]$, with $h=1/12$, $1/24$, and $1/48$, respectively. Since
each step in Euler's method requires one evaluation of $f$, the number
of evaluations of $f$ in each of these attempts is $n=12$, $24$, and
$48$, respectively. In each case we accept $y_n$ as an approximation
to $e$.  The second column of
Table~\ref{table:3.2.1} shows the results. The first column of the table
indicates
the number of evaluations of $f$ required to obtain the approximation, and
the last column contains the value of $e$ rounded to ten significant
figures.

In this section we'll study the
\href{http://www-history.mcs.st-and.ac.uk/Mathematicians/Euler.html}
{\color{blue}\itshape improved Euler method\/}, which
requires two evaluations of $f$ at each step. We've used this method
with $h=1/6$, $1/12$, and $1/24$. The required number of evaluations
of $f$ were 12, 24, and $48$, as in the three applications of Euler's
method;     however, you can see from the third column of
Table~\ref{table:3.2.1} that the approximation to $e$ obtained by the
improved Euler method with only 12 evaluations of $f$ is better than
the approximation obtained by Euler's method with 48 evaluations.

In Section~3.1 we'll study the
\href{http://www-history.mcs.st-and.ac.uk/Mathematicians/Runge.html}
{\color{blue}\itshape Runge}-\href{http://www-history.mcs.st-and.ac.uk/Mathematicians/Kutta.html}{\color{blue}\itshape Kutta method\/},
which requires four evaluations of $f$ at each step. We've used this
method with $h=1/3$, $1/6$, and $1/12$. The required number of
evaluations of $f$ were again 12, 24, and $48$, as in the three
applications of Euler's method and the improved Euler method;     however,
you can see from the fourth column of Table~\ref{table:3.2.1} that the
approximation to $e$ obtained by the Runge-Kutta method with only 12
evaluations of $f$ is better than the approximation obtained by the
improved Euler method with 48 evaluations.

\begin{newtable} \label{table:3.2.1}
{Approximations to $e$ obtained by three numerical methods.}
\begin{center}
\begin{tabular}{|c|c|c|c|c|}
\hline
\multicolumn{1}{|c|}{$n$}&
\multicolumn{1}{|c|}{Euler}&
\multicolumn{1}{|c|}{Improved Euler}&
\multicolumn{1}{|c|}{Runge-Kutta}&
\multicolumn{1}{|c|}{Exact}\\ \hline
12 & 2.613035290  &2.707188994  &2.718069764 & 2.718281828 \\
24 & 2.663731258   &2.715327371   &2.718266612 & 2.718281828 \\
48 & 2.690496599  &2.717519565  &2.718280809 & 2.718281828
\\
\hline
\end{tabular}
\end{center}
\end{newtable}


\boxit{The Improved Euler Method}


\noindent
The
\href{http://www-history.mcs.st-and.ac.uk/Mathematicians/Euler.html}
{\color{blue}\itshape improved Euler method\/} for solving the initial value
problem (\ref{eq:3.2.1}) is based on approximating the integral curve of
(\ref{eq:3.2.1}) at $(x_i,y(x_i))$ by the line through $(x_i,y(x_i))$ with
slope
\[
m_i={f(x_i,y(x_i))+f(x_{i+1},y(x_{i+1}))\over2};
\]
that is, $m_i$ is the average of the slopes of the tangents to the
integral curve at the endpoints of $[x_i,x_{i+1}]$. The equation of
the approximating line is therefore
\begin{equation} \label{eq:3.2.2}
y=y(x_i)+{f(x_i,y(x_i))+f(x_{i+1},y(x_{i+1}))\over2}(x-x_i).
\end{equation}
Setting $x=x_{i+1}=x_i+h$ in (\ref{eq:3.2.2}) yields
\begin{equation} \label{eq:3.2.3}
y_{i+1}=y(x_i)+{h\over2}\left(f(x_i,y(x_i))+f(x_{i+1},y(x_{i+1}))\right)
\end{equation}
as an approximation to $y(x_{i+1})$. As in our derivation of Euler's
method, we replace $y(x_i)$ (unknown if $i>0$) by its
approximate value $y_i$; then (\ref{eq:3.2.3}) becomes
\[
y_{i+1}=y_i+{h\over2}\left(f(x_i,y_i)+f(x_{i+1},y(x_{i+1})\right).
\]
However, this still won't work, because we don't know $y(x_{i+1})$,
which appears on the right. We overcome this by replacing $y(x_{i+1})$
by $y_i+hf(x_i,y_i)$, the value that the  Euler method would
assign to $y_{i+1}$. Thus,  the improved Euler method
starts with the known value $y(x_0)=y_0$ and computes
$y_1$, $y_2$, \dots, $y_n$ successively with the formula
\begin{equation} \label{eq:3.2.4}
y_{i+1}=y_i+{h\over2}\left(f(x_i,y_i)+f(x_{i+1},y_i+hf(x_i,y_i))\right).
\end{equation}
The computation indicated here can be conveniently organized as
follows: given $y_i$, compute
\begin{eqnarray*}
k_{1i}&=&f(x_i,y_i),\\
k_{2i}&=&f\left(x_i+h,y_i+hk_{1i}\right),\\
y_{i+1}&=&y_i+{h\over2}(k_{1i}+k_{2i}).
\end{eqnarray*}

The improved Euler method requires two evaluations of $f(x,y)$ per
step, while Euler's method requires only one. However, we'll see at
the end of this section that if $f$ satisfies appropriate assumptions,
 the local truncation error with the improved Euler method is
$O(h^3)$, rather than $O(h^2)$ as with Euler's method. Therefore
the global truncation error with the improved Euler method is
$O(h^2)$;     however, we won't prove this.

We note that the magnitude of the local truncation error in the
improved Euler method and other methods discussed in this section is
determined by the third derivative $y'''$ of the solution of the
initial value problem. Therefore the local truncation error will be
larger where $|y'''|$ is large, or smaller where $|y'''|$ is small.

The next example, which deals with the initial value problem
considered in Example~\ref{example:3.1.1}, illustrates
the computational procedure indicated in the improved Euler method.

\begin{example}\label{example:3.2.1} \rm
Use the improved Euler method with $h=0.1$ to find approximate values
of the solution of the initial value problem
\begin{equation}\label{eq:3.2.5}
y'+2y=x^3e^{-2x},\quad y(0)=1
\end{equation}
at $x=0.1,0.2,0.3$.
\end{example}

\solution
As in Example~\ref{example:3.1.1}, we rewrite (\ref{eq:3.2.5})  as
\[
y'=-2y+x^3e^{-2x},\quad y(0)=1,
\]
which is of the form (\ref{eq:3.2.1}), with
\[
f(x,y)=-2y+x^3e^{-2x},\; x_0=0,\mbox{\; and}\; y_0=1.
\]
The improved Euler method yields

\enlargethispage{.75in}
\begin{eqnarray*}
k_{10} & = & f(x_0,y_0)
  = f(0,1)=-2,\\
k_{20} & = & f(x_1,y_0+hk_{10})=f(.1,1+(.1)(-2))\\
 &=& f(.1,.8)=-2(.8)+(.1)^3e^{-.2}=-1.599181269,\\
y_1&=&y_0+{h\over2}(k_{10}+k_{20}),\\
&=&1+(.05)(-2-1.599181269)=.820040937,\\[4\jot]
k_{11} & = & f(x_1,y_1)
  = f(.1,.820040937)= -2(.820040937)+(.1)^3e^{-.2}=-1.639263142,\\
k_{21} & = & f(x_2,y_1+hk_{11})=f(.2,.820040937+.1(-1.639263142)),\\
 &=&
f(.2,.656114622)=-2(.656114622)+(.2)^3e^{-.4}=-1.306866684,\\
y_2&=&y_1+{h\over2}(k_{11}+k_{21}),\\
&=&.820040937+(.05)(-1.639263142-1.306866684)=.672734445,\\[4\jot]
k_{12} & = & f(x_2,y_2)
  = f(.2,.672734445)= -2(.672734445)+(.2)^3e^{-.4}=-1.340106330,\\
k_{22} & = & f(x_3,y_2+hk_{12})=f(.3,.672734445+.1(-1.340106330)),\\
 &=&
f(.3,.538723812)=-2(.538723812)+(.3)^3e^{-.6}=-1.062629710,\\
y_3&=&y_2+{h\over2}(k_{12}+k_{22})\\
&=&.672734445+(.05)(-1.340106330-1.062629710)=.552597643.
\end{eqnarray*}

\begin{example}\label{example:3.2.2} \rm
Table~\ref{table:3.2.2} shows results of using the improved Euler method
with step sizes $h=0.1$ and $h=0.05$ to find approximate values of the
solution of the initial value problem
\[
y'+2y=x^3e^{-2x},\quad y(0)=1
\]
at $x=0$, $0.1$, $0.2$, $0.3$, \dots, $1.0$. For comparison, it also shows
the
corresponding approximate values obtained with Euler's method in
\ref{example:3.1.2}, and the values of the exact solution
\[
y={e^{-2x}\over4}(x^4+4).
\]
The results
obtained by the improved Euler method with $h=0.1$ are better than
those obtained by Euler's method with $h=0.05$.
\end{example}

\begin{newtable} \label{table:3.2.2}
{Numerical solution of $y'+2y=x^3e^{-2x},\ y(0)=1$, by Euler's method
and the improved Euler method.}
\begin{center}
\begin{tabular}{|c|r|r|r|r|r|}
\hline
\multicolumn{1}{|c|}{$x$}&
\multicolumn{1}{|c|}{$h=0.1$}&
\multicolumn{1}{|c|}{$h=0.05$}&
\multicolumn{1}{|c|}{$h=0.1$}&
\multicolumn{1}{|c|}{$h=0.05$}&
\multicolumn{1}{|c|}{Exact}\\ \hline
0.0 & 1.000000000& 1.000000000& 1.000000000   &1.000000000 & 1.000000000 \\
0.1 & 0.800000000& 0.810005655& 0.820040937   &0.819050572 & 0.818751221 \\
0.2 & 0.640081873& 0.656266437& 0.672734445   &0.671086455 & 0.670588174 \\
0.3 & 0.512601754& 0.532290981& 0.552597643   &0.550543878 & 0.549922980 \\
0.4 & 0.411563195& 0.432887056& 0.455160637   &0.452890616 & 0.452204669 \\
0.5 & 0.332126261& 0.353785015& 0.376681251   &0.374335747 & 0.373627557 \\
0.6 & 0.270299502& 0.291404256& 0.313970920   &0.311652239 & 0.310952904 \\
0.7 & 0.222745397& 0.242707257& 0.264287611   &0.262067624 & 0.261398947 \\
0.8 & 0.186654593& 0.205105754& 0.225267702   &0.223194281 & 0.222570721 \\
0.9 & 0.159660776& 0.176396883& 0.194879501   &0.192981757 & 0.192412038 \\
1.0 & 0.139778910& 0.154715925& 0.171388070   &0.169680673 & 0.169169104\\
\hline
&\multicolumn{2}{c|}{Euler}&
\multicolumn{2}{c|}{Improved Euler}&\multicolumn{1}{c|}{Exact}\\\hline
\end{tabular}
\end{center}
\end{newtable}

\begin{example}\label{example:3.2.3} \rm
Table~\ref{table:3.2.3} shows analogous results for the nonlinear
initial value problem
\[
y'=-2y^2+xy+x^2,\ y(0)=1.
\]
We applied Euler's  method to this problem in
Example~\ref{example:3.1.3}.
\end{example}

\begin{newtable} \label{table:3.2.3}
{Numerical solution of
$y'=-2y^2+xy+x^2,\ y(0)=1$,
by Euler's method and the improved Euler method.}
\begin{center}
\begin{tabular}{|c|r|r|r|r|r|}
\hline
\multicolumn{1}{|c|}{$x$}&
\multicolumn{1}{|c|}{$h=0.1$}&
\multicolumn{1}{|c|}{$h=0.05$}&
\multicolumn{1}{|c|}{$h=0.1$}&
\multicolumn{1}{|c|}{$h=0.05$}&
\multicolumn{1}{|c|}{``Exact''}\\ \hline
0.0 & 1.000000000 & 1.000000000 &1.000000000 & 1.000000000 & 1.000000000 \\
0.1 & 0.800000000 & 0.821375000 &0.840500000 & 0.838288371 & 0.837584494 \\
0.2 & 0.681000000 & 0.707795377 &0.733430846 & 0.730556677 & 0.729641890 \\
0.3 & 0.605867800 & 0.633776590 &0.661600806 & 0.658552190 & 0.657580377 \\
0.4 & 0.559628676 & 0.587454526 &0.615961841 & 0.612884493 & 0.611901791 \\
0.5 & 0.535376972 & 0.562906169 &0.591634742 & 0.588558952 & 0.587575491 \\
0.6 & 0.529820120 & 0.557143535 &0.586006935 & 0.582927224 & 0.581942225 \\
0.7 & 0.541467455 & 0.568716935 &0.597712120 & 0.594618012 & 0.593629526 \\
0.8 & 0.569732776 & 0.596951988 &0.626008824 & 0.622898279 & 0.621907458 \\
0.9 & 0.614392311 & 0.641457729 &0.670351225 & 0.667237617 & 0.666250842 \\
1.0 & 0.675192037 & 0.701764495 &0.730069610 & 0.726985837 & 0.726015790\\
\hline
&\multicolumn{2}{c|}{Euler}&
\multicolumn{2}{c|}{Improved
Euler}&\multicolumn{1}{c|}{``Exact''}\\\hline
\end{tabular}
\end{center}
\end{newtable}

\begin{example}\label{example:3.2.4} \rm
Use step sizes $h=0.2$, $h=0.1$, and $h=0.05$ to find approximate
values of the solution of
\begin{equation} \label{eq:3.2.6}
y'-2xy=1,\quad y(0)=3
\end{equation}
at $x=0$, $0.2$, $0.4$, $0.6$, \dots, $2.0$ by
{\bf (a)} the improved Euler method;   {\bf (b)} the improved Euler
semilinear method. (We used Euler's method and the Euler semilinear
method on this problem in~\ref{example:3.1.4}.)
\end{example}

\solutionpart{a}
Rewriting (\ref{eq:3.2.6}) as
\[
y'=1+2xy,\quad y(0)=3
\]
and applying the improved Euler method with $f(x,y)=1+2xy$ yields
the results shown in Table~\ref{table:3.2.4}.


\solutionpart{b}
Since $y_1=e^{x^2}$ is a solution of the complementary equation
$y'-2xy=0$, we can  apply the improved Euler semilinear method to
(\ref{eq:3.2.6}), with
\[
y=ue^{x^2}\mbox{\quad and \quad}
u'=e^{-x^2},\quad u(0)=3.
\]
The results listed in Table~\ref{table:3.2.5} are clearly better than
those obtained by the improved Euler method.

\begin{newtable} \label{table:3.2.4}
{Numerical solution of $y'-2xy=1,\ y(0)=3$, by the improved Euler
method.}
\begin{center}
\begin{tabular}{|c|r|r|r|r|}
\hline
\multicolumn{1}{|c|}{$x$}&
\multicolumn{1}{|c|}{$h=0.2$}&
\multicolumn{1}{|c|}{$h=0.1$}&
\multicolumn{1}{|c|}{$h=0.05$}&
\multicolumn{1}{|c|}{``Exact''}\\ \hline
0.0 &   3.000000000 &   3.000000000 &   3.000000000 &   3.000000000 \\
0.2 &   3.328000000 &   3.328182400 &   3.327973600 &   3.327851973 \\
0.4 &   3.964659200 &   3.966340117 &   3.966216690 &   3.966059348 \\
0.6 &   5.057712497 &   5.065700515 &   5.066848381 &   5.067039535 \\
0.8 &   6.900088156 &   6.928648973 &   6.934862367 &   6.936700945 \\
1.0 &  10.065725534 &  10.154872547 &  10.177430736 &  10.184923955 \\
1.2 &  15.708954420 &  15.970033261 &  16.041904862 &  16.067111677 \\
1.4 &  26.244894192 &  26.991620960 &  27.210001715 &  27.289392347 \\
1.6 &  46.958915746 &  49.096125524 &  49.754131060 &  50.000377775 \\
1.8 &  89.982312641 &  96.200506218 &  98.210577385 &  98.982969504 \\
2.0 & 184.563776288 & 203.151922739 & 209.464744495 & 211.954462214 \\
\hline
\end{tabular}
\end{center}
\end{newtable}

\begin{newtable} \label{table:3.2.5}
{Numerical solution of $y'-2xy=1,\ y(0)=3$, by the improved Euler
semilinear method.}
\begin{center}
\begin{tabular}{|c|r|r|r|r|}
\hline
\multicolumn{1}{|c|}{$x$}&
\multicolumn{1}{|c|}{$h=0.2$}&
\multicolumn{1}{|c|}{$h=0.1$}&
\multicolumn{1}{|c|}{$h=0.05$}&
\multicolumn{1}{|c|}{``Exact''}\\ \hline
0.0 &   3.000000000 &  3.000000000  &   3.000000000 &   3.000000000 \\
0.2 &   3.326513400 &  3.327518315  &   3.327768620 &   3.327851973 \\
0.4 &   3.963383070 &  3.965392084  &   3.965892644 &   3.966059348 \\
0.6 &   5.063027290 &  5.066038774  &   5.066789487 &   5.067039535 \\
0.8 &   6.931355329 &  6.935366847  &   6.936367564 &   6.936700945 \\
1.0 &  10.178248417 & 10.183256733  &  10.184507253 &  10.184923955 \\
1.2 &  16.059110511 & 16.065111599  &  16.066611672 &  16.067111677 \\
1.4 &  27.280070674 & 27.287059732  &  27.288809058 &  27.289392347 \\
1.6 &  49.989741531 & 49.997712997  &  49.999711226 &  50.000377775 \\
1.8 &  98.971025420 & 98.979972988  &  98.982219722 &  98.982969504 \\
2.0 & 211.941217796 &211.951134436  & 211.953629228 & 211.954462214 \\
\hline
\end{tabular}
\end{center}
\end{newtable}

\newpage
\boxit{A Family of Methods with $O(h^3)$ Local Truncation Error}

\noindent
We'll now derive a class of methods with $O(h^3)$ local truncation
error for solving (\ref{eq:3.2.1}). For simplicity, we
assume that $f$, $f_x$, $f_y$, $f_{xx}$, $f_{yy}$, and $f_{xy}$ are
continuous and bounded for all $(x,y)$. This implies that if $y$ is
the solution of (\ref{eq:3.2.1} then $y''$ and $y'''$ are bounded
(Exercise~\ref{exer:3.2.31}).

We begin by approximating the integral curve of (\ref{eq:3.2.1}) at
$(x_i,y(x_i))$ by the line through $(x_i,y(x_i))$ with slope
\[
m_i=\sigma y'(x_i)+\rho y'(x_i+\theta h),
\]
where $\sigma$, $\rho$, and $\theta$ are constants that we'll soon
specify;     however, we insist at the outset that $0<\theta\le 1$, so
that
\[
x_i<x_i+\theta h\le x_{i+1}.
\]
The equation of the approximating line is
\begin{equation} \label{eq:3.2.7}
\begin{array}{rcl}
y&=&y(x_i)+m_i(x-x_i)\\
&=&y(x_i)+\left[\sigma y'(x_i)+\rho y'(x_i+\theta h)\right](x-x_i).
\end{array}
\end{equation}
Setting $x=x_{i+1}=x_i+h$ in (\ref{eq:3.2.7}) yields
\[
\hat y_{i+1}=y(x_i)+h\left[\sigma y'(x_i)+\rho y'(x_i+\theta h)\right]
\]
as an approximation to $y(x_{i+1})$.

To determine $\sigma$, $\rho$, and $\theta$ so that the error
\begin{equation} \label{eq:3.2.8}
\begin{array}{rcl}
E_i&=&y(x_{i+1})-\hat y_{i+1}\\
&=&y(x_{i+1})-y(x_i)-h\left[\sigma y'(x_i)+\rho y'(x_i+\theta
h)\right]
\end{array}
\end{equation}
in this approximation is $O(h^3)$, we begin by recalling from Taylor's
theorem that
\[
y(x_{i+1})=y(x_i)+hy'(x_i)+{h^2\over2}y''(x_i)+{h^3\over6}y'''(\hat
x_i),
\]
where $\hat x_i$ is in $(x_i,x_{i+1})$. Since $y'''$ is bounded this
implies that
\[
 y(x_{i+1})-y(x_i)-hy'(x_i)-{h^2\over2}y''(x_i)=O(h^3).
\]
Comparing this with (\ref{eq:3.2.8}) shows that $E_i=O(h^3)$ if
\begin{equation} \label{eq:3.2.9}
\sigma y'(x_i)+\rho y'(x_i+\theta h)=y'(x_i)+{h\over2}y''(x_i)
+O(h^2).
\end{equation}
However, applying Taylor's theorem to $y'$ shows that
\[
y'(x_i+\theta h)=y'(x_i)+\theta h y''(x_i)+{(\theta
h)^2\over2}y'''(\overline x_i),
\]
where $\overline x_i$ is in $(x_i,x_i+\theta h)$. Since $y'''$ is
bounded, this implies that
\[
y'(x_i+\theta h)=y'(x_i)+\theta h y''(x_i)+O(h^2).
\]
Substituting this into (\ref{eq:3.2.9}) and noting that the sum of
two $O(h^2)$ terms is again $O(h^2)$ shows that
$E_i=O(h^3)$ if
\[
(\sigma+\rho)y'(x_i)+\rho\theta h y''(x_i)=
y'(x_i)+{h\over2}y''(x_i),
\]
which is true if
\begin{equation} \label{eq:3.2.10}
\sigma+\rho=1 \mbox{\quad and \quad} \rho\theta={1\over2}.
\end{equation}

Since $y'=f(x,y)$, we can now conclude from (\ref{eq:3.2.8}) that
\begin{equation} \label{eq:3.2.11}
y(x_{i+1})=y(x_i)+h\left[\sigma f(x_i,y_i)+\rho f(x_i+\theta
h,y(x_i+\theta h))\right]+O(h^3)
\end{equation}
if $\sigma$, $\rho$, and $\theta$ satisfy (\ref{eq:3.2.10}). However,
this formula would not be useful even if we knew $y(x_i)$ exactly (as
we would for $i=0$), since we still wouldn't know $y(x_i+\theta h)$
exactly. To overcome this difficulty, we again use Taylor's theorem to
write
\[
y(x_i+\theta h)=y(x_i)+\theta h y'(x_i)+{h^2\over2}y''(\tilde x_i),
\]
where $\tilde x_i$ is in $(x_i,x_i+\theta h)$. Since
$y'(x_i)=f(x_i,y(x_i))$ and $y''$ is bounded, this implies that
\begin{equation} \label{eq:3.2.12}
|y(x_i+\theta h)-y(x_i)-\theta h f(x_i,y(x_i))|\le Kh^2
\end{equation}
for some constant $K$. Since $f_y$ is bounded, the mean value theorem
implies that
\[
|f(x_i+\theta h,u)-f(x_i+\theta h,v)|\le M|u-v|
\]
for some constant $M$. Letting
\[
u=y(x_i+\theta h)\mbox{\quad and \quad}
v=y(x_i)+\theta h f(x_i,y(x_i))
\]
and recalling (\ref{eq:3.2.12}) shows that
\[
f(x_i+\theta h,y(x_i+\theta h))=f(x_i+\theta h,y(x_i)+\theta h
f(x_i,y(x_i)))+O(h^2).
\]
Substituting this into (\ref{eq:3.2.11})   yields
\begin{eqnarray*}
y(x_{i+1})&=&y(x_i)+h\left[\sigma f(x_i,y(x_i))+\right.\\&&\left.\rho
f(x_i+\theta h,y(x_i)+\theta hf(x_i,y(x_i)))\right]+O(h^3).
\end{eqnarray*}
This implies that the formula
\[
y_{i+1}=y_i+h\left[\sigma f(x_i,y_i)+\rho f(x_i+\theta
h,y_i+\theta hf(x_i,y_i))\right]
\]
has $O(h^3)$ local truncation error if $\sigma$, $\rho$, and $\theta$
satisfy (\ref{eq:3.2.10}). Substituting $\sigma=1-\rho$ and
$\theta=1/2\rho$ here yields
\begin{equation} \label{eq:3.2.13}
y_{i+1}=y_i+h\left[(1-\rho)f(x_i,y_i)+\rho
f\left(x_i+{h\over2\rho},
y_i+{h\over2\rho}f(x_i,y_i)\right)\right].
\end{equation}

The computation indicated here can be conveniently organized as
follows: given $y_i$, compute
\begin{eqnarray*}
k_{1i}&=&f(x_i,y_i),\\
k_{2i}&=&f\left(x_i+{h\over2\rho},
y_i+{h\over2\rho}k_{1i}\right),\\
y_{i+1}&=&y_i+h[(1-\rho)k_{1i}+\rho k_{2i}].
\end{eqnarray*}

Consistent with our requirement that $0<\theta<1$, we require that
$\rho\ge1/2$. Letting $\rho=1/2$ in (\ref{eq:3.2.13}) yields the improved
Euler method (\ref{eq:3.2.4}). Letting $\rho=3/4$ yields
\href{http://www-history.mcs.st-and.ac.uk/Mathematicians/Heun.html}
{\color{blue}\itshape Heun's method},
\[
y_{i+1}=y_i+h\left[{1\over4}f(x_i,y_i)+{3\over4}f\left(
x_i+{2\over3}h,y_i+{2\over3}hf(x_i,y_i)\right)\right],
\]
which can be organized as
\begin{eqnarray*}
k_{1i}&=&f(x_i,y_i),\\
k_{2i}&=&f\left(x_i+{2h\over3},
y_i+{2h\over3}k_{1i}\right),\\
y_{i+1}&=&y_i+{h\over4}(k_{1i}+3k_{2i}).
\end{eqnarray*}
Letting $\rho=1$ yields the \emph{\color{blue}midpoint
method},
\[
y_{i+1}=y_i+hf\left(x_i+{h\over2},y_i+{h\over2}f(x_i,y_i)\right),
\]
which can be organized as
\begin{eqnarray*}
k_{1i}&=&f(x_i,y_i),\\
k_{2i}&=&f\left(x_i+{h\over2},
y_i+{h\over2}k_{1i}\right),\\
y_{i+1}&=&y_i+hk_{2i}.
\end{eqnarray*}

Examples involving the midpoint method and Heun's method are given in
Exercises~\ref{exer:3.2.23}-\ref{exer:3.2.30}.

\exercises
Most of the following numerical exercises involve initial value
problems considered in the exercises in Section~3.1. You'll
find it instructive to compare the results that you obtain here with
the corresponding results that you obtained in Section~3.1.

\emph{In Exercises~\ref{exer:3.2.1}--\ref{exer:3.2.5} use the improved
Euler method
to find approximate values of the solution of the given initial value
problem at the points $x_i=x_0+ih$, where $x_0$ is the point where
the initial condition is imposed and $i=1$, $2$, $3$.}

\begin{exerciselist}
\item\label{exer:3.2.1} \Cex
$y'=2x^2+3y^2-2,\quad y(2)=1;\quad h=0.05$  \,
\item\label{exer:3.2.2} \Cex
$y'=y+\sqrt{x^2+y^2},\quad y(0)=1;\quad h=0.1$  \,
\item\label{exer:3.2.3} \Cex
$y'+3y=x^2-3xy+y^2,\quad y(0)=2;\quad h=0.05$  \,
\item\label{exer:3.2.4} \Cex
$y'=\dst{1+x\over1-y^2},\quad y(2)=3;\quad h=0.1$\
\item\label{exer:3.2.5} \Cex
$y'+x^2y=\sin xy,\quad y(1)=\pi;\quad h=0.2$

\item\label{exer:3.2.6} \Cex
Use the improved Euler method with step sizes $h=0.1$, $h=0.05$, and
$h=0.025$ to find approximate values of the solution of the initial
value problem
\[
y'+3y=7e^{4x},\quad y(0)=2
\]
at $x=0$, $0.1$, $0.2$, $0.3$, \dots, $1.0$. Compare these approximate
values with the
values of the exact solution $y=e^{4x}+e^{-3x}$, which can be obtained
by the method of Section~2.1. Present your results in a table
like Table~\ref{table:3.2.2}.


\item\label{exer:3.2.7} \Cex
Use the improved Euler method with step sizes $h=0.1$, $h=0.05$, and
$h=0.025$ to find approximate values of the solution of the initial
value problem
\[
y'+{2\over x}y={3\over x^3}+1,\quad y(1)=1
\]
at $x=1.0$, $1.1$, $1.2$, $1.3$, \dots, $2.0$. Compare these approximate
values with
the values of the exact solution
\[
y={1\over3x^2}(9\ln x+x^3+2)
\]
which can be obtained by the method of Section~2.1. Present
your results in a table like Table~\ref{table:3.2.2}.

\item\label{exer:3.2.8} \Cex
Use the improved Euler method with step sizes $h=0.05$, $h=0.025$, and
$h=0.0125$ to find approximate values of the solution of the initial
value problem
\[
y'={y^2+xy-x^2\over x^2},\quad y(1)=2,
\]
at $x=1.0$, $1.05$, $1.10$, $1.15$, \dots, $1.5$. Compare these approximate
values
with the values of the exact solution
\[
y={x(1+x^2/3)\over1-x^2/3}
\]
obtained in Example~\ref{example:2.4.3}. Present your results in a table
like Table~\ref{table:3.2.2}.

\item\label{exer:3.2.9} \Cex
In Example~\ref{example:3.2.2} it was shown that
\[
y^5+y=x^2+x-4
\]
is an implicit solution of  the initial value problem
\[
y'={2x+1\over5y^4+1},\quad y(2)=1.
\eqno{\rm(A)}
\]
Use the improved Euler method with step sizes $h=0.1$, $h=0.05$, and
$h=0.025$ to find approximate values of the solution of (A) at
$x=2.0$, $2.1$, $2.2$, $2.3$, \dots, $3.0$. Present your results in tabular
form. To
check the error in these approximate values, construct another table
of values of the residual
\[
R(x,y)=y^5+y-x^2-x+4
\]
for each value of $(x,y)$ appearing in the first table.

\item\label{exer:3.2.10} \Cex
You can see from Example~\ref{example:2.5.1} that
\[
x^4y^3+x^2y^5+2xy=4
\]
is an implicit solution of  the initial value problem
\[
y'=-{4x^3y^3+2xy^5+2y\over3x^4y^2+5x^2y^4+2x},\quad y(1)=1.
\eqno{\rm(A)}
\]
Use the improved Euler method with step sizes $h=0.1$, $h=0.05$, and
$h=0.025$ to find approximate values of the solution of (A) at
$x=1.0$, $1.14$, $1.2$, $1.3$, \dots, $2.0$. Present your results in
tabular form. To
check the error in these approximate values, construct another table
of values of the residual
\[
R(x,y)=x^4y^3+x^2y^5+2xy-4
\]
for each value of $(x,y)$ appearing in the first table.


\item\label{exer:3.2.11} \Cex
Use the improved Euler method with step sizes $h=0.1$, $h=0.05$, and
$h=0.025$ to find approximate values of the solution of the initial
value problem
\[
(3y^2+4y)y'+2x+\cos x=0, \quad y(0)=1 \mbox{\
(Exercise~2.2.~\hspace*{-3pt}\ref{exer:2.2.13})}
\]
at $x=0$, $0.1$, $0.2$, $0.3$, \dots, $1.0$.

\item\label{exer:3.2.12} \Cex
Use the improved Euler method with step sizes $h=0.1$, $h=0.05$, and
$h=0.025$ to find approximate values of the solution of the initial
value problem
\[
y'+{(y+1)(y-1)(y-2)\over x+1}=0, \quad y(1)=0
\mbox{\  (Exercise~2.2.~\hspace*{-3pt}\ref{exer:2.2.14})}
\]
at $x=1.0$, $1.1$, $1.2$, $1.3$, \dots, $2.0$.

\item\label{exer:3.2.13} \Cex
Use the improved Euler method and the improved Euler semilinear method
with step sizes $h=0.1$, $h=0.05$, and $h=0.025$ to find approximate
values of the solution of the initial value problem
\[
y'+3y=e^{-3x}(1-2x),\quad y(0)=2,
\]
at $x=0$, $0.1$, $0.2$, $0.3$, \dots, $1.0$. Compare these approximate
values with the
values of the exact solution $y=e^{-3x}(2+x-x^2)$, which can be
obtained by the method of Section~2.1. Do you notice
anything special about the results? Explain.

\exercisetext{The linear initial value problems in
Exercises~\ref{exer:3.2.14}--\ref{exer:3.2.19} can't be solved exactly in
terms of known elementary functions. In each exercise use the improved
Euler and improved Euler semilinear methods with the indicated step
sizes to find approximate values of the solution of the given initial
value problem at 11 equally spaced points (including the
endpoints) in the interval.}

\item\label{exer:3.2.14} \Cex
$y'-2y=\dst{1\over1+x^2},\quad y(2)=2$;  \quad
$h=0.1,0.05,0.025$ on $[2,3]$

\item\label{exer:3.2.15} \Cex
$y'+2xy=x^2,\quad y(0)=3$;\quad
$h=0.2,0.1,0.05$ on $[0,2]$
\quad (Exercise~2.1.~\hspace*{-3pt}\ref{exer:2.1.38})

\item\label{exer:3.2.16} \Cex
$\dst{y'+{1\over x}y={\sin x\over x^2},\quad y(1)=2}$, \quad
$h=0.2,0.1,0.05$ on $[1,3]$
\quad (Exercise~2.1.\hspace*{-3pt}\ref{exer:2.1.39})

\item\label{exer:3.2.17} \Cex
$\dst{y'+y={e^{-x}\tan x\over x},\quad y(1)=0}$;\quad
$h=0.05,0.025,0.0125$ on $[1,1.5]$\;
(Exercise~2.1.~\hspace*{-3pt}\ref{exer:2.1.40}),

\item\label{exer:3.2.18} \Cex
$\dst{y'+{2x\over 1+x^2}y={e^x\over (1+x^2)^2}, \quad y(0)=1}$;\;
$h=0.2,0.1,0.05$ on $[0,2]$\quad
 (Exercise~2.1.~\hspace*{-3pt}\ref{exer:2.1.41})

\item\label{exer:3.2.19} \Cex
$xy'+(x+1)y=e^{x^2},\quad y(1)=2$;\;
$h=0.05,0.025,0.0125$ on $[1,1.5]$\;
(Exercise~2.1.~\hspace*{-3pt}\ref{exer:2.1.42})

\exercisetext{In Exercises~\ref{exer:3.2.20}--\ref{exer:3.2.22} use the
improved Euler method and the improved Euler semilinear method with
the indicated step sizes to find approximate values of the solution of
the given initial value problem at 11 equally spaced points
(including the endpoints) in the interval.}

\item\label{exer:3.2.20} \Cex
$y'+3y=xy^2(y+1),\quad y(0)=1$;\quad
$h=0.1,0.05,0.025$ on $[0,1]$

\item\label{exer:3.2.21} \Cex
$\dst{y'-4y={x\over y^2(y+1)},\quad y(0)=1}$;\quad
 $h=0.1,0.05,0.025$ on $[0,1]$

\item\label{exer:3.2.22} \Cex
$\dst{y'+2y={x^2\over1+y^2},\quad y(2)=1}$;\quad
 $h=0.1,0.05,0.025$ on $[2,3]$


\item\label{exer:3.2.23} \Cex
Do Exercise~\ref{exer:3.2.7} with ``improved Euler method''
replaced by ``midpoint method.''

\item\label{exer:3.2.24} \Cex
Do Exercise~\ref{exer:3.2.7} with ``improved Euler method''
replaced by ``Heun's method.''

\item\label{exer:3.2.25} \Cex
Do Exercise~\ref{exer:3.2.8} with ``improved Euler method''
replaced by ``midpoint method.''

\item\label{exer:3.2.26} \Cex
Do Exercise~\ref{exer:3.2.8} with ``improved Euler method''
replaced by ``Heun's method.''

\item\label{exer:3.2.27} \Cex
Do Exercise~\ref{exer:3.2.11} with ``improved Euler method''
replaced by ``midpoint method.''

\item\label{exer:3.2.28} \Cex
Do Exercise~\ref{exer:3.2.11} with ``improved Euler method''
replaced by ``Heun's method.''

\item\label{exer:3.2.29} \Cex
Do Exercise~\ref{exer:3.2.12} with ``improved Euler method''
replaced by ``midpoint method.''

\item\label{exer:3.2.30} \Cex
Do Exercise~\ref{exer:3.2.12} with ``improved Euler method''
replaced by ``Heun's method.''

\item\label{exer:3.2.31}
Show that if $f$, $f_x$, $f_y$, $f_{xx}$, $f_{yy}$, and $f_{xy}$ are
continuous and bounded for all $(x,y)$ and $y$ is the solution of the
initial value problem
\[
y'=f(x,y),\quad y(x_0)=y_0,
\]
then $y''$ and $y'''$ are bounded.

\item\label{exer:3.2.32}
{\sc Numerical Quadrature} (see Exercise~3.1.~\hspace*{-3pt}\ref{exer:3.1.23}).
\begin{alist}
\item % (a)
Derive the quadrature formula
\[
\int_a^bf(x)\,dx\approx .5h(f(a)+f(b))+
h\sum_{i=1}^{n-1}f(a+ih)\mbox{\quad(where $h=(b-a)/n)$}
\eqno{\rm(A)}
\]
by applying the improved Euler method to the initial value problem
\[
y'=f(x),\quad y(a)=0.
\]
\item % (b)
The quadrature formula (A) is  called \emph{\color{blue}the trapezoid
rule}. Draw a figure that justifies this terminology.
\item % (c)
\Lex
For several choices of $a$, $b$, $A$,  and $B$,
 apply (A) to $f(x)=A+Bx$, with $n = 10,20,40,80,160,320$.
Compare your results with the exact answers
and explain what you find.
\item % (d)
\Lex
For several choices of $a$, $b$, $A$, $B$, and $C$,
 apply (A) to $f(x)=A+Bx+Cx^2$, with $n=10$, $20$, $40$, $80$, $160$,
$320$.
 Compare your results with the exact answers
and explain what you find.
\end{alist}

\end{exerciselist}
\setcounter{table}{0}
\section{The Runge-Kutta Method}

In general, if $k$ is any  positive  integer  and $f$
satisfies appropriate assumptions, there are numerical methods with
local truncation error  $O(h^{k+1})$ for solving an initial value
problem
\begin{equation} \label{eq:3.3.1}
y'=f(x,y),\quad y(x_0)=y_0.
\end{equation}
Moreover, it can be shown that a method with local truncation error
$O(h^{k+1})$ has global truncation error $O(h^k)$. In
Sections~3.1\ and 3.2 we studied numerical
methods where
$k=1$ and $k=2$. We'll skip methods for which $k=3$ and proceed to the
\href{http://www-history.mcs.st-and.ac.uk/Mathematicians/Runge.html}
{\color{blue}\itshape Runge}-\href{http://www-history.mcs.st-and.ac.uk/Mathematicians/Kutta.html}{\itshape\color{blue}Kutta\/} method, the most widely used
method, for which $k=4$. The magnitude of the local truncation error is
determined by the fifth derivative $y^{(5)}$ of the solution of the
initial value problem. Therefore the local truncation error will be
larger where $|y^{(5)}|$ is large, or smaller where $|y^{(5)}|$ is
small.
\enlargethispage{1in}
The Runge-Kutta method computes approximate values
$y_1$, $y_2$, \dots, $y_n$ of the solution of \eqref{eq:3.3.1}
at $x_0$, $x_0+h$, \dots, $x_0+nh$ as follows: Given $y_i$,
compute
\begin{eqnarray*} k_{1i}&=&f(x_i,y_i),\\
k_{2i}&=&f\left(x_i+{h\over2},y_i+{h\over2}k_{1i}\right),\\
k_{3i}&=&f\left(x_i+{h\over2},y_i+{h\over2}k_{2i}\right),\\
k_{4i}&=&f(x_i+h,y_i+hk_{3i}),
\end{eqnarray*}
and
\[
y_{i+1}=y_i+{h\over6}(k_{1i}+2k_{2i}+2k_{3i}+k_{4i}).
\]

 The next example, which deals with the initial value problem
considered in Examples~\ref{example:3.1.1} and \ref{example:3.2.1},
illustrates
the computational procedure indicated in the Runge-Kutta method.

\begin{example}\label{example:3.3.1} \rm
Use the Runge-Kutta method with $h=0.1$ to find approximate values for
the solution of the initial value problem
\begin{equation}\label{eq:3.3.2}
y'+2y=x^3e^{-2x},\quad y(0)=1,
\end{equation}
at $x=0.1,0.2$.
\end{example}

\solution
Again we rewrite \eqref{eq:3.3.2}  as
\[
y'=-2y+x^3e^{-2x},\quad y(0)=1,
\]
which is of the form \eqref{eq:3.3.1}, with
\[
f(x,y)=-2y+x^3e^{-2x},\ x_0=0,\mbox{\ and}\ y_0=1.
\]
\enlargethispage{.1in}

The Runge-Kutta method yields
\begin{eqnarray*}
k_{10} & = & f(x_0,y_0)
  = f(0,1)=-2,\\
k_{20} & = & f(x_0+h/2,y_0+hk_{10}/2)=f(.05,1+(.05)(-2))\\
 &=& f(.05,.9)=-2(.9)+(.05)^3e^{-.1}=-1.799886895,\\
k_{30} & = & f(x_0+h/2,y_0+hk_{20}/2)=f(.05,1+(.05)(-1.799886895))\\
 &=& f(.05,.910005655)=-2(.910005655)+(.05)^3e^{-.1}=-1.819898206,\\
k_{40} & = & f(x_0+h,y_0+hk_{30})=f(.1,1+(.1)(-1.819898206))\\
&=&f(.1,.818010179)=-2(.818010179)+(.1)^3e^{-.2}=-1.635201628,\\
y_1&=&y_0+{h\over6}(k_{10}+2k_{20}+2k_{30}+k_{40}),\\
&=&1+{.1\over6}(-2+2(-1.799886895)+2(-1.819898206)
-1.635201628)=.818753803,\\[4\jot]
k_{11} & = & f(x_1,y_1)
  = f(.1,.818753803)=-2(.818753803))+(.1)^3e^{-.2}=-1.636688875,\\
k_{21} & = & f(x_1+h/2,y_1+hk_{11}/2)=f(.15,.818753803+(.05)(-1.636688875))\\
 &=& f(.15,.736919359)=-2(.736919359)+(.15)^3e^{-.3}=-1.471338457,\\
k_{31} & = & f(x_1+h/2,y_1+hk_{21}/2)=f(.15,.818753803+(.05)(-1.471338457))\\
 &=& f(.15,.745186880)=-2(.745186880)+(.15)^3e^{-.3}=-1.487873498,\\
k_{41} & = &
f(x_1+h,y_1+hk_{31})=f(.2,.818753803+(.1)(-1.487873498))\\
&=&f(.2,.669966453)=-2(.669966453)+(.2)^3e^{-.4}=-1.334570346,\\
y_2&=&y_1+{h\over6}(k_{11}+2k_{21}+2k_{31}+k_{41}),\\
&=&.818753803+{.1\over6}(-1.636688875+2(-1.471338457)+2(-1.487873498)-1.334570346)
\\&=&.670592417.
\end{eqnarray*}

The Runge-Kutta method is sufficiently accurate for most
applications.

\begin{example}\label{example:3.3.2} \rm
Table~\ref{table:3.3.1} shows results of using the Runge-Kutta method
with step sizes $h=0.1$ and $h=0.05$ to find approximate values of
the solution of the initial value problem
\[
y'+2y=x^3e^{-2x},\quad y(0)=1
\]
at $x=0$, $0.1$, $0.2$, $0.3$, \dots, $1.0$. For comparison, it also shows
the
corresponding approximate values obtained with the improved Euler
method in Example~\ref{example:3.2.2}, and the values of the exact
solution
\[
y={e^{-2x}\over4}(x^4+4).
\]
The results obtained by the Runge-Kutta method are clearly better
than those obtained by the improved Euler method   in fact; the results
obtained by the Runge-Kutta method with $h=0.1$ are better than those
obtained by the improved Euler method with $h=0.05$.
\end{example}



\begin{newtable} \label{table:3.3.1}
{Numerical solution of
$y'+2y=x^3e^{-2x},\ y(0)=1$, by
the Runge-Kuttta method and the improved Euler
method.}
\begin{center}
\begin{tabular}{|c|r|r|r|r|r|}
\hline
\multicolumn{1}{|c|}{$x$}&
\multicolumn{1}{|c|}{$h=0.1$}&
\multicolumn{1}{|c|}{$h=0.05$}&
\multicolumn{1}{|c|}{$h=0.1$}&
\multicolumn{1}{|c|}{$h=0.05$}&
\multicolumn{1}{|c|}{Exact}\\ \hline
0.0 & 1.000000000   &1.000000000 &1.000000000 & 1.000000000& 1.000000000 \\
0.1 & 0.820040937   &0.819050572 &0.818753803 & 0.818751370& 0.818751221 \\
0.2 & 0.672734445   &0.671086455 &0.670592417 & 0.670588418& 0.670588174 \\
0.3 & 0.552597643   &0.550543878 &0.549928221 & 0.549923281& 0.549922980 \\
0.4 & 0.455160637   &0.452890616 &0.452210430 & 0.452205001& 0.452204669 \\
0.5 & 0.376681251   &0.374335747 &0.373633492 & 0.373627899& 0.373627557 \\
0.6 & 0.313970920   &0.311652239 &0.310958768 & 0.310953242& 0.310952904 \\
0.7 & 0.264287611   &0.262067624 &0.261404568 & 0.261399270& 0.261398947 \\
0.8 & 0.225267702   &0.223194281 &0.222575989 & 0.222571024& 0.222570721 \\
0.9 & 0.194879501   &0.192981757 &0.192416882 & 0.192412317& 0.192412038 \\
1.0 & 0.171388070   &0.169680673 &0.169173489 & 0.169169356& 0.169169104\\
\hline
&\multicolumn{2}{c|}{Improved Euler}&
\multicolumn{2}{c|}{Runge-Kutta}&\multicolumn{1}{c|}{Exact}\\\hline
\end{tabular}
\end{center}
\end{newtable}

\begin{example}\label{example:3.3.3} \rm
Table~\ref{table:3.3.2}   shows analogous results
for the nonlinear initial value problem
\[
y'=-2y^2+xy+x^2,\ y(0)=1.
\]
We applied the improved Euler  method to this problem in
Example~\ref{exer:3.2.3}.
\end{example}

\begin{newtable} \label{table:3.3.2}
{Numerical solution of
$y'=-2y^2+xy+x^2,\ y(0)=1$, by
 the Runge-Kuttta method and the improved Euler
method.}
\begin{center}
\begin{tabular}{|c|r|r|r|r|r|}
\hline
\multicolumn{1}{|c|}{$x$}&
\multicolumn{1}{|c|}{$h=0.1$}&
\multicolumn{1}{|c|}{$h=0.05$}&
\multicolumn{1}{|c|}{$h=0.1$}&
\multicolumn{1}{|c|}{$h=0.05$}&
\multicolumn{1}{|c|}{``Exact''}\\ \hline
0.0  &1.000000000 & 1.000000000 &1.000000000 & 1.000000000& 1.000000000 \\
0.1  &0.840500000 & 0.838288371 &0.837587192 & 0.837584759& 0.837584494 \\
0.2  &0.733430846 & 0.730556677 &0.729644487 & 0.729642155& 0.729641890 \\
0.3  &0.661600806 & 0.658552190 &0.657582449 & 0.657580598& 0.657580377 \\
0.4  &0.615961841 & 0.612884493 &0.611903380 & 0.611901969& 0.611901791 \\
0.5  &0.591634742 & 0.588558952 &0.587576716 & 0.587575635& 0.587575491 \\
0.6  &0.586006935 & 0.582927224 &0.581943210 & 0.581942342& 0.581942225 \\
0.7  &0.597712120 & 0.594618012 &0.593630403 & 0.593629627& 0.593629526 \\
0.8  &0.626008824 & 0.622898279 &0.621908378 & 0.621907553& 0.621907458 \\
0.9  &0.670351225 & 0.667237617 &0.666251988 & 0.666250942& 0.666250842 \\
1.0  &0.730069610 & 0.726985837 &0.726017378 & 0.726015908& 0.726015790\\
\hline
&\multicolumn{2}{c|}{Improved Euler}&
\multicolumn{2}{c|}{Runge-Kutta}
&\multicolumn{1}{c|}{``Exact''}\\\hline
\end{tabular}
\end{center}
\end{newtable}


\begin{example}\label{example:3.3.4} \rm
Tables~\ref{table:3.3.3} and \ref{table:3.3.4} show results obtained by
applying the Runge-Kutta and Runge-Kutta semilinear methods to
to the initial value problem
\[
y'-2xy=1,\ y(0)=3,
\]
which we considered in Examples~\ref{example:3.1.4}  and
\ref{example:3.2.4}.
\end{example}

\begin{newtable} \label{table:3.3.3}
{Numerical solution of
$y'-2xy=1,\ y(0)=3$, by
 the Runge-Kutta method.}
\begin{center}
\begin{tabular}{|c|r|r|r|r|}
\hline
\multicolumn{1}{|c|}{$x$}&
\multicolumn{1}{|c|}{$h=0.2$}&
\multicolumn{1}{|c|}{$h=0.1$}&
\multicolumn{1}{|c|}{$h=0.05$}&
\multicolumn{1}{|c|}{``Exact''}\\ \hline
0.0 &   3.000000000 &   3.000000000 &   3.000000000 &   3.000000000 \\
0.2 &   3.327846400 &   3.327851633 &   3.327851952 &   3.327851973 \\
0.4 &   3.966044973 &   3.966058535 &   3.966059300 &   3.966059348 \\
0.6 &   5.066996754 &   5.067037123 &   5.067039396 &   5.067039535 \\
0.8 &   6.936534178 &   6.936690679 &   6.936700320 &   6.936700945 \\
1.0 &  10.184232252 &  10.184877733 &  10.184920997 &  10.184923955 \\
1.2 &  16.064344805 &  16.066915583 &  16.067098699 &  16.067111677 \\
1.4 &  27.278771833 &  27.288605217 &  27.289338955 &  27.289392347 \\
1.6 &  49.960553660 &  49.997313966 &  50.000165744 &  50.000377775 \\
1.8 &  98.834337815 &  98.971146146 &  98.982136702 &  98.982969504 \\
2.0 & 211.393800152 & 211.908445283 & 211.951167637 & 211.954462214 \\
\hline
\end{tabular}
\end{center}
\end{newtable}

\begin{newtable} \label{table:3.3.4}
{Numerical solution of
$y'-2xy=1,\ y(0)=3$, by
 the Runge-Kutta semilinear
method.}
\begin{center}
\begin{tabular}{|c|r|r|r|r|}
\hline
\multicolumn{1}{|c|}{$x$}&
\multicolumn{1}{|c|}{$h=0.2$}&
\multicolumn{1}{|c|}{$h=0.1$}&
\multicolumn{1}{|c|}{$h=0.05$}&
\multicolumn{1}{|c|}{``Exact''}\\ \hline
0.0 &   3.000000000 &   3.000000000 &   3.000000000 &   3.000000000 \\
0.2 &   3.327853286 &   3.327852055 &   3.327851978 &   3.327851973 \\
0.4 &   3.966061755 &   3.966059497 &   3.966059357 &   3.966059348 \\
0.6 &   5.067042602 &   5.067039725 &   5.067039547 &   5.067039535 \\
0.8 &   6.936704019 &   6.936701137 &   6.936700957 &   6.936700945 \\
1.0 &  10.184926171 &  10.184924093 &  10.184923963 &  10.184923955 \\
1.2 &  16.067111961 &  16.067111696 &  16.067111678 &  16.067111677 \\
1.4 &  27.289389418 &  27.289392167 &  27.289392335 &  27.289392347 \\
1.6 &  50.000370152 &  50.000377302 &  50.000377745 &  50.000377775 \\
1.8 &  98.982955511 &  98.982968633 &  98.982969450 &  98.982969504 \\
2.0 & 211.954439983 & 211.954460825 & 211.954462127 & 211.954462214 \\
\hline
\end{tabular}
\end{center}
\end{newtable}

\boxit{The Case Where $x_0$ Isn't  The Left Endpoint}


\enlargethispage{.5in}
\noindent
So far in this chapter we've considered numerical methods
 for solving an initial value problem
\begin{equation} \label{eq:3.3.3}
y'=f(x,y),\quad y(x_0)=y_0
\end{equation}
on an interval $[x_0,b]$, for which $x_0$ is the left endpoint. We
haven't discussed numerical methods for solving \eqref{eq:3.3.3}
on an interval $[a,x_0]$, for which $x_0$ is the
right endpoint. To be specific, how can we obtain approximate values
$y_{-1}$, $y_{-2}$, \dots, $y_{-n}$ of the solution of \eqref{eq:3.3.3} at
$x_0-h, \dots,x_0-nh$,
where $h=(x_0-a)/n$?  Here's the answer to this question:

\medskip
\color{blue}
 Consider the initial value problem
\begin{equation} \label{eq:3.3.4}
z'=-f(-x,z),\quad z(-x_0)=y_0,
\end{equation}
on the interval $[-x_0,-a]$, for which $-x_0$ is the left endpoint.
Use a numerical method to obtain approximate values
$z_1$, $z_2$, \dots, $z_n$ of the solution of $\eqref{eq:3.3.4}$ at
$-x_0+h$, $-x_0+2h$, \dots, $-x_0+nh=-a$. Then
$y_{-1}=z_1$, $y_{-2}=z_2$, $\dots$, $y_{-n}=z_n$
are approximate values of the solution of $\eqref{eq:3.3.3}$ at
$x_0-h$, $x_0-2h$, \dots, $x_0-nh=a$.
\color{black}
\medskip

The justification for this answer is sketched in
Exercise~\ref{exer:3.3.23}. Note how easy it is to make the change
the given problem \eqref{eq:3.3.3} to the modified problem
\eqref{eq:3.3.4}: first replace $f$ by $-f$ and then replace $x$,
$x_0$, and $y$ by $-x$, $-x_0$, and $z$, respectively.


\begin{example}\label{example:3.3.5} \rm
Use the Runge-Kutta  method with step size $h=0.1$
 to find approximate values of the solution of
\begin{equation} \label{eq:3.3.5}
(y-1)^2y'=2x+3,\quad y(1)=4
\end{equation}
at $x=0$, $0.1$, $0.2$, \dots, $1$.
\end{example}

\solution We first rewrite \eqref{eq:3.3.5} in the form \eqref{eq:3.3.3}
as
\begin{equation} \label{eq:3.3.6}
y'={2x+3\over(y-1)^2},\quad y(1)=4.
\end{equation}
Since the initial condition $y(1)=4$  is imposed at the right endpoint
of the interval $[0,1]$, we apply the Runge-Kutta method to the
initial value problem
\begin{equation} \label{eq:3.3.7}
z'={2x-3\over(z-1)^2},\quad z(-1)=4
\end{equation}
on the interval $[-1,0]$.
(You should verify that \eqref{eq:3.3.7} is related to \eqref{eq:3.3.6} as
\eqref{eq:3.3.4} is related to \eqref{eq:3.3.3}.)
 Table~\ref{table:3.3.5} shows the results.  Reversing the
order of the
rows in Table~\ref{table:3.3.5} and changing the signs of the values of $x$
yields the first two columns of Table~\ref{table:3.3.6}.   The last
column of Table~\ref{table:3.3.6} shows the exact values of $y$, which
are given by
\[
y=1+(3x^2+9x+15)^{1/3}.
\]
(Since the differential equation in \eqref{eq:3.3.6} is separable,
this formula can be obtained by the method of Section~2.2.)

\begin{newtable} \label{table:3.3.5}
{Numerical solution of $z'=\dst{2x-3\over(z-1)^2},\ z(-1)=4$, on
$[-1,0]$.}
\begin{center}
\begin{tabular}{|r|r|}
\hline
\multicolumn{1}{|c|}{$x$}&
\multicolumn{1}{|c|}{$z$}\\ \hline
-1.0 & 4.000000000  \\
-0.9 & 3.944536474  \\
-0.8 & 3.889298649  \\
-0.7 & 3.834355648  \\
-0.6 & 3.779786399  \\
-0.5 & 3.725680888  \\
-0.4 & 3.672141529  \\
-0.3 & 3.619284615  \\
-0.2 & 3.567241862  \\
-0.1 & 3.516161955  \\
 0.0 & 3.466212070 \\
\hline
\end{tabular}
\end{center}
\end{newtable}

\begin{newtable} \label{table:3.3.6}
{Numerical solution of $(y-1)^2y'=2x+3,\ y(1)=4$, on
$[0,1]$.}
\begin{center}
\begin{tabular}{|r|r|r|}
\hline
\multicolumn{1}{|c|}{$x$}&
\multicolumn{1}{|c|}{$y$}&
\multicolumn{1}{|c|}{Exact}\\ \hline
0.00  &   3.466212070 &  3.466212074\\
0.10  &   3.516161955 &  3.516161958\\
0.20  &   3.567241862 &  3.567241864\\
0.30  &   3.619284615 &  3.619284617\\
0.40  &   3.672141529 &  3.672141530\\
0.50  &   3.725680888 &  3.725680889\\
0.60  &   3.779786399 &  3.779786399\\
0.70  &   3.834355648 &  3.834355648\\
0.80  &   3.889298649 &  3.889298649\\
0.90  &   3.944536474 &  3.944536474\\
1.00  &   4.000000000 &  4.000000000\\
\hline
\end{tabular}
\end{center}
\end{newtable}

We leave it to you to develop a procedure for handling the
numerical solution of \eqref{eq:3.3.3} on an interval $[a,b]$
such that $a<x_0<b$ (Exercises~\ref{exer:3.3.26} and \ref{exer:3.3.27}).

\exercises
Most of the following numerical exercises involve initial value
problems considered in the exercises in Sections~3.2.
 You'll find it instructive to compare the results that you
obtain here with the corresponding results that you obtained in those
sections.

\noindent
\emph{In Exercises~\ref{exer:3.3.1}--\ref{exer:3.3.5} use the Runge-Kutta
method to
find approximate values of the solution of the given initial value
problem at the points $x_i=x_0+ih,$ where $x_0$ is the point where
the initial condition is imposed and $i=1$, $2$.}


\begin{exerciselist}
\item\label{exer:3.3.1} \Cex
$y'=2x^2+3y^2-2,\quad y(2)=1;\quad h=0.05$
\item\label{exer:3.3.2} \Cex
$y'=y+\sqrt{x^2+y^2},\quad y(0)=1;\quad h=0.1$
\item\label{exer:3.3.3} \Cex
$y'+3y=x^2-3xy+y^2,\quad y(0)=2;\quad h=0.05$
\item\label{exer:3.3.4} \Cex
$y'=\dst{1+x\over1-y^2},\quad y(2)=3;\quad h=0.1$
\item\label{exer:3.3.5} \Cex
$y'+x^2y=\sin xy,\quad y(1)=\pi;\quad h=0.2$

\item\label{exer:3.3.6} \Cex
Use  the Runge-Kutta method with step sizes $h=0.1$, $h=0.05$, and
$h=0.025$ to
find approximate values of the solution of the initial value problem
\[
y'+3y=7e^{4x},\quad y(0)=2,
\]
at $x=0$, $0.1$, $0.2$, $0.3$, \dots, $1.0$. Compare these approximate
values with the
values of the exact solution $y=e^{4x}+e^{-3x}$, which can be obtained
by the method of Section~2.1. Present your results in a table
like Table~\ref{table:3.3.1}.


\item\label{exer:3.3.7} \Cex
Use the Runge-Kutta method with step sizes $h=0.1$, $h=0.05$, and
$h=0.025$ to find approximate values of the solution of the initial
value problem
\[
y'+{2\over x}y={3\over x^3}+1,\quad y(1)=1
\]
at $x=1.0$, $1.1$, $1.2$, $1.3$, \dots, $2.0$. Compare these approximate
values
with the values of the exact solution
\[
y={1\over3x^2}(9\ln x+x^3+2),
\]
which can be obtained by the method of Section~2.1. Present
your results in a table like Table~\ref{table:3.3.1}.

\item\label{exer:3.3.8} \Cex
Use the Runge-Kutta method with step sizes $h=0.05$, $h=0.025$, and
$h=0.0125$ to find approximate values of the solution of the initial
value problem
\[
y'={y^2+xy-x^2\over x^2},\quad y(1)=2
\]
at $x=1.0$, $1.05$, $1.10$, $1.15$ \dots, $1.5$. Compare these approximate
values
with the values of the exact solution
\[
y={x(1+x^2/3)\over1-x^2/3},
\]
which was obtained in Example~\ref{example:2.2.3}. Present your results in
a table like Table~\ref{table:3.3.1}.

\item\label{exer:3.3.9} \Cex
In Example~\ref{example:2.2.3} it was shown that
\[
y^5+y=x^2+x-4
\]
is an implicit solution of  the initial value problem
\[
y'={2x+1\over5y^4+1},\quad y(2)=1.
\eqno{\rm(A)}
\]
Use the Runge-Kutta method with step sizes $h=0.1$, $h=0.05$, and
$h=0.025$ to find approximate values of the solution of (A) at
$x=2.0$, $2.1$, $2.2$, $2.3$, \dots, $3.0$. Present your results in tabular
form.
To check the error in these approximate values, construct another
table of values of the residual
\[
  R(x,y)=y^5+y-x^2-x+4
\]
for each value of $(x,y)$ appearing in the first table.

\item\label{exer:3.3.10} \Cex
You can see from Example~\ref{example:2.5.1} that
\[
x^4y^3+x^2y^5+2xy=4
\]
is an implicit solution of  the initial value problem
\[
y'=-{4x^3y^3+2xy^5+2y\over3x^4y^2+5x^2y^4+2x},\quad y(1)=1.
\eqno{\rm(A)}
\]
Use the Runge-Kutta method with step sizes $h=0.1$, $h=0.05$, and
$h=0.025$ to find approximate values of the solution of (A) at
$x=1.0$, $1.1$, $1.2$, $1.3$, \dots, $2.0$. Present your results in tabular
form.
To check the error in these approximate values, construct another
table of values of the residual
\[
R(x,y)=x^4y^3+x^2y^5+2xy-4
\]
for each value of $(x,y)$ appearing in the first table.

\item\label{exer:3.3.11} \Cex
Use the Runge-Kutta method with step sizes $h=0.1$, $h=0.05$, and
$h=0.025$ to find approximate values of the solution of the initial
value problem
\[
(3y^2+4y)y'+2x+\cos x=0, \quad y(0)=1 \mbox{\
(Exercise~2.2.~\hspace*{-3pt}\ref{exer:2.2.13})},
\]
at $x=0$, $0.1$, $0.2$, $0.3$, \dots, $1.0$.

\item\label{exer:3.3.12} \Cex
Use the Runge-Kutta method with step sizes $h=0.1$, $h=0.05$, and
$h=0.025$ to find approximate values of the solution of the initial
value problem
\[
y'+{(y+1)(y-1)(y-2)\over x+1}=0, \quad y(1)=0
\mbox{\ (Exercise~2.2.~\hspace*{-3pt}\ref{exer:2.2.14})},
\]
at $x=1.0$, $1.1$, $1.2$, $1.3$, \dots, $2.0$.

\item\label{exer:3.3.13} \Cex
Use the Runge-Kutta method and the Runge-Kutta  semilinear method
with step sizes
$h=0.1$, $h=0.05$, and $h=0.025$ to find approximate values of the
solution of the initial value problem
\[
y'+3y=e^{-3x}(1-4x+3x^2-4x^3),\quad y(0)=-3
\]
at $x=0$, $0.1$, $0.2$, $0.3$, \dots, $1.0$. Compare these approximate
values with
the values of the exact solution $y=-e^{-3x}(3-x+2x^2-x^3+x^4)$, which
can be obtained by the method of Section~2.1. Do
you notice anything special about the results? Explain.

\exercisetext{
The linear initial value problems in
Exercises~\ref{exer:3.3.14}--\ref{exer:3.3.19} can't be solved exactly in
terms of known elementary functions. In each exercise use the
Runge-Kutta and the Runge-Kutta semilinear methods with the
indicated step sizes to find approximate values of the solution of the
given initial value problem at $11$ equally spaced points (including
the endpoints) in the interval.}

\item\label{exer:3.3.14} \Cex
$y'-2y=\dst{1\over1+x^2},\quad y(2)=2$;  \,
$h=0.1,0.05,0.025$ on $[2,3]$

\item\label{exer:3.3.15} \Cex
$y'+2xy=x^2,\quad y(0)=3$;\quad
$h=0.2,0.1,0.05$ on $[0,2]$
(Exercise~2.1.~\hspace*{-3pt}\ref{exer:2.1.38})

\item\label{exer:3.3.16} \Cex
$\dst{y'+{1\over x}y={\sin x\over x^2},\quad y(1)=2;}$\quad
$h=0.2,0.1,0.05$ on $[1,3]$
(Exercise~2.1.~\hspace*{-3pt}\ref{exer:2.1.39})

\item\label{exer:3.3.17} \Cex
$\dst{y'+y={e^{-x}\tan x\over x},\quad y(1)=0;}$\quad
$h=0.05,0.025,0.0125$ on $[1,1.5]$
(Exercise~2.1.~\hspace*{-3pt}\ref{exer:2.1.40})

\item\label{exer:3.3.18} \Cex
$\dst{y'+{2x\over 1+x^2}y={e^x\over (1+x^2)^2}, \quad y(0)=1};$ \quad
$h=0.2,0.1,0.05$ on $[0,2]$
 (Exercise~2.1,~\hspace*{-3pt}\ref{exer:2.1.41})

\item\label{exer:3.3.19} \Cex
$xy'+(x+1)y=e^{x^2},\quad y(1)=2$;\quad
$h=0.05,0.025,0.0125$ on $[1,1.5]$
(Exercise~2.1.~\hspace*{-3pt}\ref{exer:2.1.42})

\exercisetext{
In Exercises~\ref{exer:3.3.20}--\ref{exer:3.3.22} use the Runge-Kutta method
and the Runge-Kutta semilinear method with the indicated step sizes
to find approximate values of the solution of the given initial value
problem at $11$ equally spaced points (including the endpoints) in
the interval.}

\item\label{exer:3.3.20} \Cex
$y'+3y=xy^2(y+1),\quad y(0)=1$;  \,
$h=0.1,0.05,0.025$ on $[0,1]$


\item\label{exer:3.3.21} \Cex
$\dst{y'-4y={x\over y^2(y+1)},\quad y(0)=1}$;  \,
 $h=0.1,0.05,0.025$ on $[0,1]$

\item\label{exer:3.3.22} \Cex
$\dst{y'+2y={x^2\over1+y^2},\quad y(2)=1}$;  \,
 $h=0.1,0.05,0.025$ on $[2,3]$

\item\label{exer:3.3.23} \Cex
Suppose $a<x_0$, so that $-x_0<-a$.
Use the chain rule to show that if  $z$ is a solution of
\[
z'=-f(-x,z),\quad z(-x_0)=y_0,
\]
on $[-x_0,-a]$, then $y=z(-x)$ is a solution of
\[
y'=f(x,y),\quad y(x_0)=y_0,
\]
on $[a,x_0]$.

\item\label{exer:3.3.24} \Cex
Use the Runge-Kutta  method with step sizes $h=0.1$, $h=0.05$,
and $h=0.025$ to find approximate values of the solution of
\[
y'={y^2+xy-x^2\over x^2},\quad y(2)=-1
\]
at $x=1.1$, $1.2$, $1.3$, \dots $2.0$.
 Compare these approximate values
with the values of the exact solution
\[
y={x(4-3x^2)\over4+3x^2},
\]
which can be obtained by referring to Example~\ref{example:2.4.3}.

\item\label{exer:3.3.25} \Cex
Use the Runge-Kutta  method with step sizes $h=0.1$,  $h=0.05$,
and $h=0.025$
 to find approximate values of the solution of
\[
y'=-x^2y-xy^2,\quad y(1)=1
\]
at $x=0$, $0.1$, $0.2$, \dots, $1$.

\item\label{exer:3.3.26} \Cex
Use the Runge-Kutta  method with step sizes $h=0.1$,  $h=0.05$,
and $h=0.025$
 to find approximate values of the solution of
\[
y'+{1\over x}y={7\over x^2}+3,\quad y(1)={3\over2}
\]
at $x=0.5$, $0.6$,\dots, $1.5$.
 Compare these approximate values
with the values of the exact solution
\[
y={7\ln x\over x}+{3x\over2},
\]
which can be obtained by the method discussed in
Section~2.1.

\item\label{exer:3.3.27} \Cex
Use the Runge-Kutta  method with step sizes $h=0.1$,  $h=0.05$,
and $h=0.025$
 to find approximate values of the solution of
\[
xy'+2y=8x^2,\quad y(2)=5
\]
at $x=1.0$, $1.1$, $1.2$, \dots, $3.0$.
 Compare these approximate values
with the values of the exact solution
\[
y=2x^2-{12\over x^2},
\]
which can be obtained by the method discussed in Section~2.1.


\item\label{exer:3.3.28}
{\sc Numerical Quadrature} (see
Exercise~3.1.~\hspace*{-3pt}\ref{exer:3.1.23}).
\begin{alist}
\item % (a)
Derive the quadrature formula
\[
\int_a^bf(x)\,dx\approx {h\over6}(f(a)+f(b))+
{h\over3}\sum_{i=1}^{n-1}f(a+ih)+{2h\over3}\sum_{i=1}^n
f\left(a+(2i-1)h/2\right)
\eqno{\rm(A)}
\]
(where $h=(b-a)/n)$
by applying the Runge-Kutta  method to the initial value problem
\[
y'=f(x),\quad y(a)=0.
\]
This quadrature formula is  called
\href{http://www-history.mcs.st-and.ac.uk/Mathematicians/Simpson.html}
{\color{blue}\itshape Simpson's Rule\/}.
\item % (b)
\Lex
For several choices of $a$, $b$, $A$, $B$, $C$, and $D$
 apply (A) to $f(x)=A+Bx+Cx+Dx^3$, with $n = 10$, $20$, $40$,
$80$, $160$, $320$.
Compare your results with the exact answers
and explain what you find.
\item % (d)
\Lex
For several choices of $a$, $b$, $A$, $B$, $C$, $D$, and $E$
 apply (A) to $f(x)=A+Bx+Cx^2+Dx^3+Ex^4$, with
$n=10,20,40,80,160,320$.
 Compare your results with the exact answers
and explain what you find.
\end{alist}

\end{exerciselist}
